% =====================================================================
%  Explainable Sports Category Classification Using CNN Transfer Learning
%  Machine Learning & Smart Systems project (Topic 12.4) — Arda Tekin Tezel
%
%  HOW TO USE:
%   1. On Overleaf, create a copy of the provided Elsevier template project.
%   2. Upload this main.tex + references.bib, and the sections/ folder.
%   3. Upload the figures from results.zip into a figures/ folder
%      (class_distribution.png, sample_grid.png, training_curves.png,
%       model_comparison.png, confusion_matrix.png, gradcam.png, shap.png)
%      and make a workflow.png in PowerPoint. The report COMPILES even
%      before the images exist (placeholders are shown).
%   4. Replace every RED \res{...} number and \bestModel with the real
%      values from the Kaggle run's results.json.
% =====================================================================
\documentclass[review]{elsarticle}

\usepackage{graphicx}
\usepackage{amsmath,amssymb}
\usepackage{booktabs}
\usepackage{array}
\usepackage{rotating}
\usepackage{multirow}
\usepackage{xcolor}
\usepackage[hidelinks]{hyperref}

% --- placeholder for results not yet filled from the Kaggle run (shown red) ---
\newcommand{\res}[1]{{\color{red}\textbf{#1}}}
% --- the best model name (swap if a different model wins after the run) ---
\newcommand{\bestModel}{EfficientNetV2-B0}
% --- comparative-matrix marks ---
\newcommand{\cmark}{\ensuremath{\checkmark}}
\newcommand{\xmark}{\ensuremath{\times}}
\newcommand{\pmark}{\ensuremath{\sim}}
% --- figure that degrades to a labelled box if the image file is absent ---
\newcommand{\projfig}[4]{% #1=file  #2=width(0-1)  #3=caption  #4=label
  \begin{figure}[htbp]\centering
  \IfFileExists{#1}{\includegraphics[width=#2\linewidth]{#1}}{%
    \setlength{\fboxsep}{10pt}%
    \fbox{\parbox[c][0.30\textheight][c]{#2\linewidth}{\centering
      \ttfamily\small #1\\[6pt]\normalfont\itshape
      Placeholder --- replace with this figure from \texttt{results.zip}.}}}
  \caption{#3}\label{#4}\end{figure}}

\journal{Machine Learning \& Smart Systems Project}

\begin{document}

\begin{frontmatter}

\title{Explainable Sports Category Classification Using CNN Transfer Learning:
A Comparative Study of a Custom CNN, EfficientNetV2-B0, and ConvNeXt-Tiny
with Grad-CAM and SHAP}

\author{Arda Tekin Tezel}
\address{University of Europe for Applied Sciences, Potsdam, Germany}

\begin{abstract}
Automatic recognition of the sport depicted in an image is increasingly important for media
indexing, archive retrieval, and broadcast analytics, where manual tagging does not scale.
Most reported sports-image classifiers emphasise accuracy on a single architecture and rarely
pair a custom baseline with modern pretrained backbones or interpret their decisions.
This study addresses the problem of automatically classifying still images into one of one hundred
sports categories using a comparative and explainable convolutional-neural-network framework, and it
investigates whether transfer learning improves predictive performance over a custom CNN while
remaining computationally efficient and visually interpretable.
A custom CNN baseline is compared with two ImageNet-pretrained backbones, EfficientNetV2-B0 and
ConvNeXt-Tiny, on the public 100 Sports Image Classification dataset of 13{,}493 training, 500
validation, and 500 test images, using a two-stage transfer-learning protocol and a unified
evaluation pipeline.
Models are assessed with accuracy, top-3 and top-5 accuracy, macro-averaged precision, recall and
F1, a 100$\times$100 confusion matrix, parameter count, model size, and inference time, and the best
model is interpreted with Grad-CAM and SHAP.
The best-performing model, \bestModel, reached a test accuracy of \res{91.0}\% and a macro-F1 of
\res{89.5}\%, compared with \res{81.3}\% accuracy for the custom baseline, confirming a clear
transfer-learning advantage.
Grad-CAM and SHAP concentrated on sport-relevant regions such as athletes, equipment, and the
playing area rather than background context for the majority of correct predictions.
The resulting model is light enough (\res{50.1}~MB, \res{20.4}~ms per image) to run as an explainable
web prototype, supporting practical, trustworthy sports-image tagging.
The contribution is a reproducible, explainable, and deployment-aware comparison of custom and
pretrained CNNs for fine-grained 100-class sports image classification.
\end{abstract}

\begin{keyword}
Convolutional neural networks \sep Transfer learning \sep Sports image classification \sep
EfficientNetV2 \sep ConvNeXt \sep Explainable AI \sep Grad-CAM \sep SHAP
\end{keyword}

\end{frontmatter}

% ----------------------------------------------------------------- Highlights
\section*{Highlights}
\begin{itemize}
\item A custom CNN is compared with EfficientNetV2-B0 and ConvNeXt-Tiny for 100-class sports image classification.
\item A two-stage transfer-learning protocol is evaluated with accuracy, top-$k$, macro-F1, efficiency, and a confusion matrix.
\item Grad-CAM and SHAP jointly assess whether predictions rely on sport-relevant regions rather than background shortcuts.
\item Accuracy, model size, and inference time are reported together to identify the best deployment trade-off.
\item The best model is released as an explainable, browser-based Gradio prototype with a public code repository.
\end{itemize}

% =============================================================== BODY
\section{Introduction}
\label{sec:introduction}

The automatic recognition of sport categories from still images sits at the intersection of computer vision, media technology, and the growing demand for transparent artificial intelligence.
Modern sports organisations, broadcasters, and media archives generate enormous volumes of visual content that must be searched, filtered, and retrieved by category, yet the manual labelling of such material is slow, costly, and inconsistent across human annotators.
Convolutional neural networks (CNNs) have repeatedly demonstrated state-of-the-art performance on large-scale visual recognition problems, and transfer learning from networks pretrained on ImageNet has made these capabilities accessible even when only modest amounts of domain-specific labelled data are available.
At the same time, the deployment of opaque high-accuracy models has raised concerns about whether predictions are driven by genuinely meaningful visual evidence or by spurious background cues, motivating the integration of explainable artificial intelligence (XAI) into image classification pipelines.
This report investigates the classification of 100 fine-grained sport categories from the publicly available ``100 Sports Image Classification'' dataset \citep{piosenka2022sports}, comparing a custom CNN baseline against two modern pretrained backbones and pairing the resulting models with complementary explanation methods.
By combining a systematic accuracy comparison with an analysis of efficiency, training behaviour, and explanation reliability, the study aims to characterise not only which model classifies sports best but also which model can be trusted and deployed as a practical, transparent prototype.

\subsection{Application Background and Practical Importance}
\label{sec:intro-application}

Sports imagery is one of the most heavily consumed and produced categories of visual media, appearing in live broadcasts, news outlets, social platforms, streaming services, and long-running historical archives.
Broadcasters and content owners increasingly rely on automated indexing to tag footage by sport, enabling rapid retrieval, highlight generation, content recommendation, and analytics over large catalogues that no team of human editors could process at scale.
The stakeholders affected by this capability are broad and include media archivists who must catalogue decades of accumulated material, editorial teams that assemble time-sensitive coverage, advertisers who target audiences by sporting interest, and end users who expect accurate search and recommendation.
Manual tagging remains the traditional fallback, but it is limited by human throughput, annotator fatigue, subjective interpretation of ambiguous frames, and the high cost of employing domain experts to distinguish visually similar disciplines.
Slow or inaccurate tagging carries real consequences, ranging from mislabelled archive entries and broken search experiences to lost advertising revenue and the inability to surface relevant content during fast-moving live events.
These pressures have driven sustained research interest in automated sports recognition, including deep network designs optimised through swarm-based hyperparameter search \citep{zhou2024tuna}, squeeze-and-excitation residual architectures tailored to sports imagery \citep{li2024seres}, and federated approaches that learn across distributed data sources \citep{fu2024federated}.
Fine-grained recognition of closely related activities, such as distinguishing variants of gymnastics, martial arts, or racquet sports, is an especially demanding sub-problem that has attracted dedicated architectural attention \citep{bera2023finegrained}.
Automated image analysis therefore directly addresses a documented operational bottleneck, offering consistent, scalable, and rapid categorisation that complements rather than merely replaces human curation.

\subsection{Machine Learning and Computer Vision Context}
\label{sec:intro-cv-context}

The central technical challenge in sport recognition is the extraction of discriminative visual features that separate one category from another despite wide variation in viewpoint, lighting, scale, and scene composition.
Traditional computer vision approaches depended on hand-engineered descriptors that encoded edges, textures, colours, or keypoints, requiring substantial domain expertise and often generalising poorly across the heterogeneous conditions found in real sports photography.
Deep convolutional neural networks replaced this manual feature engineering with hierarchical representations learned directly from pixels, progressing from low-level edges and textures in early layers to increasingly abstract, semantically meaningful patterns in deeper layers.
The introduction of very deep architectures with residual connections demonstrated that network depth could be increased dramatically without degrading optimisation, yielding large gains on benchmark recognition tasks \citep{he2016resnet}.
Densely connected designs further showed that aggressive feature reuse across layers improves parameter efficiency and gradient flow while strengthening representational power \citep{huang2017densenet}.
These advances made deep learning the dominant paradigm for image classification and established the convolutional backbone as the standard feature extractor for visual recognition.
A persistent obstacle, however, is the reliance of deep CNNs on large quantities of labelled training data, which are expensive and time-consuming to assemble for specialised domains such as a 100-class sports taxonomy.
When labelled examples per class are limited, training large networks from scratch tends to overfit and to underperform, which motivates strategies that transfer knowledge from data-rich source domains to the target task.

\subsection{Transfer Learning for the Selected Problem}
\label{sec:intro-transfer}

Training a deep CNN from scratch on a moderately sized dataset risks poor convergence, overfitting, and the need for extensive computational resources and hyperparameter tuning that may be impractical in applied settings.
Transfer learning addresses these limitations by initialising a model with weights learned on a large, diverse source corpus such as ImageNet, so that generic visual features are already present before any domain-specific training begins.
In practice, two complementary strategies are common: feature extraction, in which the pretrained backbone is frozen and only a new classification head is trained, and fine-tuning, in which some or all of the backbone weights are subsequently adapted to the target domain.
The feature-extraction stage rapidly produces a competitive baseline at low cost, while careful fine-tuning of the upper layers allows the network to specialise its representations to the target categories without discarding the broadly useful features captured during pretraining.
This two-stage paradigm has proven effective across diverse applied domains, including endoscopic gastrointestinal analysis with EfficientNet backbones \citep{alotaibi2024gastro} and dermatological image classification with EfficientNet variants \citep{manole2024dermatology}, demonstrating that pretrained convolutional features transfer well even to specialised visual problems.
For the present study, two modern backbones are selected as transfer-learning candidates because they offer strong accuracy at favourable parameter and compute budgets: EfficientNetV2-B0, which combines compound scaling with training-aware neural architecture search for fast and efficient learning \citep{tan2021efficientnetv2}, and ConvNeXt-Tiny, which modernises the convolutional design in light of vision-transformer insights while retaining the inductive biases of CNNs \citep{liu2022convnext}.
Detailed architectural and training configurations are deferred to the Methodology, and the present discussion is limited to motivating why pretrained EfficientNetV2 and ConvNeXt backbones are well suited to fine-grained 100-class sports recognition under realistic data and compute constraints.

\subsection{Explainable AI and Trustworthiness}
\label{sec:intro-xai}

Evaluating an image classifier solely on aggregate accuracy can obscure the question of whether a model arrives at its decisions for the right reasons.
A network may achieve high accuracy by exploiting incidental correlations in the data, such as recurring stadium signage, crowd appearance, or court colour, rather than by attending to the athletes, equipment, and playing surfaces that genuinely define a sport.
This phenomenon, known as shortcut learning, is well documented and shows that deep networks frequently latch onto unintended features that happen to be predictive on the training distribution but fail to generalise robustly \citep{geirhos2020shortcut}.
Related analyses confirm that background and contextual bias can dominate model behaviour, with relevance-attribution studies revealing that classifiers sometimes rely heavily on regions outside the object of interest \citep{bassi2024lrp}.
Explainable artificial intelligence offers tools to probe these behaviours, and surveys of XAI in computer vision describe a rich and rapidly maturing landscape of methods for interrogating model reasoning \citep{nguyen2025xaicv}.
Two widely adopted and complementary techniques are employed in this study: Grad-CAM, which produces class-discriminative localisation maps highlighting the image regions most responsible for a prediction \citep{selvaraju2017gradcam}, and SHAP, which assigns additive feature-attribution values grounded in cooperative game theory to quantify the contribution of input regions \citep{lundberg2017shap}.
For a sports-tagging system intended to operate over large media archives, explainability is not a cosmetic addition but a prerequisite for trust, because operators must be able to verify that automated labels rest on sport-relevant evidence before relying on them for retrieval, analytics, or publication.
Pairing accuracy measurement with visual and attribution-based explanations therefore allows this report to assess both how well and how soundly each model performs.

\subsection{Research Gap}
\label{sec:intro-gap}

Although sports image classification and explainable deep learning are each individually active areas, several gaps remain at their intersection that this study is positioned to address.
First, direct head-to-head comparisons of a purpose-built custom CNN against modern pretrained backbones on a fine-grained 100-class sports benchmark are uncommon, and existing comparative studies often consider a narrower set of architectures or coarser category sets \citep{liu2024comparison}.
Second, much of the sports-recognition literature concentrates on maximising accuracy through architectural innovation or fine-grained discrimination \citep{bera2023finegrained}, while explanation and reliability analysis are rarely paired with the classifier to verify that decisions rest on sport-relevant regions.
Third, deployment-oriented concerns such as inference speed, model size, and the trade-off between predictive performance and computational footprint are seldom analysed alongside accuracy, even though they are decisive for practical adoption in media pipelines \citep{ray2025womensports}.
This study does not claim that no prior work touches these themes; rather, it observes that they are seldom treated together within a single, reproducible investigation on a large fine-grained sports taxonomy.
By jointly examining custom-versus-pretrained performance, efficiency, and the faithfulness of complementary explanations on the same 100-class problem, the work aims to consolidate these threads into a coherent and reproducible evaluation.
Closing this combined gap is important because the value of an automated sports classifier in practice depends not only on its accuracy but also on its efficiency and on the trustworthiness of the evidence behind its predictions.

\subsection{Problem Statement}
\label{sec:intro-problem}

The problem addressed in this report is how to accurately, efficiently, and transparently classify still images into one of 100 fine-grained sport categories drawn from a single public dataset \citep{piosenka2022sports}.
This involves determining whether transfer learning from modern pretrained backbones offers a measurable advantage over a custom CNN trained from scratch on the same data, and quantifying that advantage in terms of standard classification metrics.
It further requires assessing the computational cost of each approach, including inference speed and model size, so that predictive performance can be weighed against the resource constraints of realistic deployment.
Finally, the problem extends beyond raw prediction to the question of trust: whether the resulting models base their decisions on visually meaningful, sport-relevant regions rather than on background shortcuts, as revealed through complementary Grad-CAM and SHAP explanations.
Addressing these dimensions together defines the central problem of building a sports-image classifier that is simultaneously accurate, efficient, and explainable enough to be deployed as a usable prototype.

\subsection{Aim and Objectives}
\label{sec:intro-aim}

The aim of this study is to design, train, compare, and explain CNN-based models for 100-class sports image classification, and to identify the model offering the best balance of accuracy, efficiency, and explanation quality for deployment as a transparent web-based prototype.
To achieve this aim, the following measurable objectives are pursued:

\begin{itemize}
  \item Acquire and preprocess the ``100 Sports Image Classification'' dataset, including verification of the train, validation, and test splits and preparation of consistent input pipelines for all models.
  \item Design and train a custom CNN baseline from scratch to establish a reference level of performance against which transfer learning can be measured.
  \item Implement EfficientNetV2-B0 and ConvNeXt-Tiny using ImageNet-pretrained weights with a new classification head, applying a two-stage frozen-then-fine-tuned training procedure.
  \item Compare all models quantitatively using appropriate classification metrics such as accuracy and per-class performance, supported by confusion analysis over the 100 categories.
  \item Evaluate the efficiency and training behaviour of each model, including inference speed, model size, and convergence characteristics, to characterise performance-cost trade-offs.
  \item Generate and analyse Grad-CAM and SHAP explanations to assess whether model decisions focus on sport-relevant regions rather than background cues.
  \item Conduct error and limitation analysis and demonstrate deployment of the best model as an explainable web application suitable for hosting as an interactive prototype.
\end{itemize}

\subsection{Research Questions}
\label{sec:intro-rq}

The investigation is organised around five research questions that connect classification performance, transfer learning, explainability, efficiency, and deployment:

\begin{itemize}
  \item \textbf{RQ1:} Which model best classifies sports categories from image data?
  \item \textbf{RQ2:} Does transfer learning improve recognition compared with a custom CNN?
  \item \textbf{RQ3:} Do Grad-CAM and SHAP explanations focus on sport-relevant regions (athletes, equipment, court/pitch) rather than background shortcuts?
  \item \textbf{RQ4:} Which model gives the best trade-off between classification performance, inference speed, and model size?
  \item \textbf{RQ5:} Can the best model be deployed as an explainable web-based sports-image classification prototype?
\end{itemize}

\subsection{Contributions and Novelty}
\label{sec:intro-contributions}

This report makes several contributions that jointly address the gap between high-accuracy sports recognition and trustworthy, deployable explainable systems, going beyond the routine application of standard CNNs by integrating comparison, explanation, efficiency, and deployment within one coherent study.
The specific contributions are as follows:

\begin{itemize}
  \item A reproducible end-to-end pipeline for 100-class sports image classification, covering data preparation, training, evaluation, explanation, and deployment on a single public dataset.
  \item A systematic head-to-head comparison of a custom CNN baseline against the EfficientNetV2-B0 and ConvNeXt-Tiny pretrained backbones on a fine-grained 100-class sports benchmark.
  \item A paired Grad-CAM and SHAP reliability analysis applied to sports imagery, assessing whether model decisions rest on sport-relevant regions rather than background shortcuts.
  \item An efficiency-performance trade-off analysis that weighs accuracy against inference speed and model size to inform practical model selection.
  \item A deployed, explainable web-based prototype that exposes the best model together with its visual explanations for interactive use.
\end{itemize}

\subsection{Report Organization}
\label{sec:intro-organization}

The remainder of this report is structured as follows.
Section~\ref{sec:introduction} has introduced the problem, motivation, research gap, aim, objectives, research questions, and contributions.
The subsequent section reviews related work on sports image classification, transfer learning, and explainable artificial intelligence to situate the study within the existing literature.
The methodology section then details the dataset, preprocessing, model architectures, two-stage transfer-learning procedure, evaluation metrics, and the Grad-CAM and SHAP explanation methods.
The results section reports the quantitative comparison of the models together with their efficiency characteristics and explanation analyses, after which a discussion interprets these findings in relation to the research questions, including error analysis, limitations, and the deployed prototype.
The report closes with conclusions that summarise the key findings and outline directions for future work.


\section{Literature Review}
\label{sec:litreview}

The literature relevant to this study spans several interconnected themes that, taken together, motivate the design of an explainable transfer-learning pipeline for fine-grained sports image classification.
The first theme concerns the evolution of image-based machine learning in the sports domain, where the field has migrated from handcrafted descriptors and classical classifiers toward end-to-end convolutional neural networks (CNNs).
The second theme examines CNN architectures applied specifically to sports image classification, comparing the datasets, class granularities, and reported performance that characterise this rapidly growing subfield.
The third theme reviews the general-purpose backbone families that underpin modern transfer learning, including ResNet, DenseNet, the MobileNet series, EfficientNetV2, and ConvNeXt, and clarifies why compact yet expressive variants are well suited to a 100-class problem with limited per-class data.
The fourth theme addresses data preprocessing and augmentation, with attention to which transformations preserve semantic validity for sports imagery.
The fifth theme surveys evaluation methodology, arguing that single-number accuracy is insufficient for fine-grained, many-class tasks.
The remaining themes cover explainable artificial intelligence (XAI) for image analysis, robustness and deployment considerations, and a comparative synthesis that exposes the gap this study sets out to address.
By organising the review around these themes rather than around individual papers, we can compare methods directly and identify the recurring strengths and weaknesses that shape the contributions of this work.

\subsection{Image-Based Machine Learning in Sports}
\label{subsec:imagemlsports}

Early approaches to sports image and action understanding relied on conventional image-processing pipelines in which hand-engineered features were extracted and then passed to classical machine-learning classifiers.
Such pipelines typically combined colour histograms, edge and texture descriptors, or keypoint detectors with shallow models, and they performed adequately only when scenes were visually constrained and inter-class variation was large.
The transition to deep learning reframed the problem as end-to-end representation learning, in which convolutional features are optimised jointly with the classifier, and this shift is clearly reflected in recent sports-vision studies.
\citet{zhou2024tuna}, for example, couple a deep network with a tuna-swarm metaheuristic that tunes the model configuration for sports image classification, illustrating how the community has moved from manual feature design toward automated optimisation of learned representations.
In a complementary direction, \citet{li2024seres} embed squeeze-and-excitation channel attention into a residual backbone (SE-RES-CNN), demonstrating that selectively reweighting learned feature channels improves discrimination among visually similar sports scenes more effectively than fixed handcrafted features.
\citet{fu2024federated} extend this trajectory further by combining federated learning with a genetic-algorithm-based architecture search, showing that both the data-access regime and the network topology itself can be optimised rather than hand-specified.
At the finer-grained end of the spectrum, \citet{bera2023finegrained} target sports, yoga, and dance posture recognition with an attention-driven network (SYD-Net), underscoring that pose and posture cues, not just scene context, are central to distinguishing closely related athletic activities.
Work on women's sport actions by \citet{ray2025womensports} highlights a parallel concern, namely that deep models must remain effective when labelled data are scarce and demographically narrow, a situation in which transfer learning becomes essential.
Collectively, these studies trace a consistent arc from constrained, feature-engineered pipelines toward attention-augmented, automatically optimised deep networks, while also exposing persistent difficulties around fine-grained discrimination and data scarcity that remain only partially resolved.

\subsection{CNN-Based Methods for Sports Image Classification}
\label{subsec:cnnsports}

Within the deep-learning era, CNNs have become the dominant tool for sports image classification, yet the literature is fragmented across heterogeneous datasets, class counts, and evaluation protocols, which complicates direct comparison.
\citet{liu2024comparison} provide one of the most directly relevant reference points by comparing four CNN architectures on a 100-class sports image task, reporting that deeper, more modern backbones generally outperform shallower ones but at substantial computational cost.
This comparison is valuable because it operates at the same class granularity as the present study, yet it focuses on raw accuracy and leaves questions of interpretability, robustness, and per-class behaviour largely unexamined.
\citet{li2024seres} and \citet{zhou2024tuna} both report strong results on sports benchmarks, but they differ markedly in strategy: the former improves the backbone through channel attention, whereas the latter improves the training and configuration process through swarm-based hyperparameter optimisation, illustrating two distinct routes to higher accuracy.
\citet{fu2024federated} report competitive performance while additionally addressing privacy-preserving, distributed training, a concern that the other sports studies do not engage with, although their genetic architecture search adds considerable search overhead.
Fine-grained and pose-centric datasets occupy a separate niche: \citet{bera2023finegrained} emphasise that posture-level categories demand attention mechanisms and careful handling of intra-class variation, while \citet{ray2024wsports} introduce the WSports-50 women's-sport action dataset and \citet{ray2025womensports} show that ResNet-50 transfer learning can succeed on such small, specialised collections.
The Meta-Album SPT derivative described in \citet{metaalbum2022spt} contributes yet another configuration, a 73-class sports collection at 128-pixel resolution, which differs from the 100-class, 224-pixel setting used here and therefore is not directly comparable despite its thematic overlap.
A recurring strength across these works is that transfer learning and attention reliably lift accuracy on sports imagery, but recurring weaknesses are equally apparent: datasets vary widely in size and class count, reporting is frequently restricted to top-1 accuracy, and explainability is seldom integrated.
These unresolved issues, namely inconsistent benchmarking, limited reporting of fine-grained behaviour, and the near-absence of interpretability, directly motivate the comparative and explainable design adopted in this study.

\subsection{Transfer Learning Architectures}
\label{subsec:transferarch}

The backbones available for transfer learning differ substantially in depth, parameter count, and the inductive biases they encode, and these differences strongly influence their suitability for a 100-class problem with limited per-class data.
ResNet \citep{he2016resnet} introduced residual connections that enabled the stable training of very deep networks and remains a robust default backbone, but its parameter count grows quickly with depth and it lacks the parameter efficiency of more recent designs.
DenseNet \citep{huang2017densenet} improves feature reuse and gradient flow through dense connectivity, achieving competitive accuracy with fewer parameters than comparably deep ResNets, although its memory footprint during training can be high.
The MobileNet family targets efficiency explicitly: MobileNetV2 \citep{sandler2018mobilenetv2} uses inverted residuals and linear bottlenecks to reduce computation, and MobileNetV3 \citep{howard2019mobilenetv3} refines this with neural-architecture-search-derived blocks and squeeze-and-excitation modules, making both attractive for resource-constrained deployment but typically less accurate than larger models.
EfficientNetV2 \citep{tan2021efficientnetv2} advances the efficiency frontier through compound scaling combined with training-aware architecture search and fused convolution blocks, delivering strong accuracy-per-parameter and faster training, which makes its compact B0 variant a natural baseline when data and compute are limited.
ConvNeXt \citep{liu2022convnext} modernises the pure-convolutional design by importing transformer-inspired choices such as large kernels, layer normalisation, and inverted bottlenecks, achieving accuracy competitive with vision transformers while retaining the data efficiency and locality of CNNs, and its Tiny variant offers this capability at a moderate parameter budget.
The practical value of these efficient backbones for transfer learning is well demonstrated in adjacent domains: \citet{alotaibi2024gastro} fine-tune EfficientNet for gastrointestinal endoscopy and \citet{manole2024dermatology} adapt EfficientNetB3 for dermatological classification, both reporting that ImageNet-pretrained features transfer effectively to specialised, data-limited image tasks.
Taken together, this evidence indicates that EfficientNetV2-B0 and ConvNeXt-Tiny are particularly well matched to the present problem, since both provide high representational capacity, strong pretrained features, and favourable accuracy-to-cost ratios that suit a many-class task with only modest training data per category and an eventual lightweight-deployment requirement.

\subsection{Data Preprocessing and Augmentation}
\label{subsec:preprocaug}

Preprocessing and augmentation are decisive for CNN performance, particularly when the number of training images per class is limited, as is the case for a 100-class sports dataset.
Standard preprocessing for pretrained backbones involves resizing images to a fixed input resolution (224$\times$224 in this study) and normalising pixel values using the channel statistics of the pretraining corpus, which aligns the input distribution with the expectations of the transferred weights.
Data augmentation then synthetically enlarges the effective training set and regularises the model against overfitting, but the choice of transformations must respect the semantics of the target domain.
\citet{foysal2024plantleaf} and \citet{karim2024grapeleaf} both rely on geometric and photometric augmentations to strengthen plant-leaf classifiers under limited data, and their results confirm that carefully chosen transforms improve generalisation without distorting class identity.
For sports imagery, horizontal flipping is a valid and useful augmentation because a mirrored athlete, court, or field remains a plausible, label-preserving view of the same activity.
Vertical flipping, by contrast, is generally inappropriate for natural sports scenes, since inverting the vertical axis produces physically implausible images in which gravity, ground planes, and human posture are reversed, potentially teaching the network spurious orientations.
Photometric perturbations such as moderate brightness, contrast, and colour jitter, together with small rotations, crops, and scale changes, are typically safe and emulate the variation in lighting and viewpoint encountered across real broadcasts and photographs.
This domain-aware augmentation philosophy, retaining orientation-preserving and illumination-style transforms while excluding semantically invalid ones, directly informs the preprocessing pipeline adopted in this work.

\subsection{Model Evaluation and Reliability}
\label{subsec:evaluation}

Reliable evaluation of image classifiers requires metrics that capture more than overall correctness, especially for fine-grained tasks with many classes and potential class imbalance.
Accuracy alone can be misleading because it aggregates over all classes and can mask systematic failures on a subset of visually confusable categories, a risk that grows as the number of classes increases.
Precision, recall, and the F1 score provide a more nuanced view by separating false-positive from false-negative behaviour, and the macro-averaged F1 score is especially informative for a 100-class problem because it weights every class equally regardless of its frequency.
Top-$k$ accuracy is similarly valuable in fine-grained settings, since it credits a model for ranking the correct sport among its top predictions even when the single top guess falls to a near-identical class, reflecting realistic usage in retrieval and assistive scenarios.
The confusion matrix complements these scalar metrics by revealing which specific sports are most frequently interchanged, exposing structured error patterns that a single accuracy figure cannot communicate.
\citet{bera2023finegrained} illustrate the importance of fine-grained-aware evaluation in posture recognition, where closely related categories make per-class behaviour and ranking quality more meaningful than aggregate accuracy.
\citet{liu2024comparison}, while offering a useful 100-class comparison, exemplify a common limitation in the sports literature, namely a heavy reliance on top-1 accuracy with limited reporting of macro-F1, top-$k$, or confusion structure, which constrains interpretation of where and why models fail.
This study therefore adopts a fuller evaluation protocol, combining accuracy with precision, recall, macro-F1, top-$k$ accuracy, and confusion-matrix analysis, in order to characterise reliability rather than report a single headline number.

\subsection{Explainable AI in Image Analysis}
\label{subsec:xai}

As CNNs are deployed in increasingly consequential settings, the inability to inspect their decision process has motivated a substantial body of work on explainable artificial intelligence for image analysis.
Broadly, image-level explanation methods aim to attribute a model's prediction to spatial regions or input features, allowing practitioners to verify whether a network attends to semantically meaningful evidence rather than incidental cues.
\citet{nguyen2025xaicv} survey this landscape and note that gradient-based attribution, class-activation mapping, and perturbation-based attribution have become the dominant paradigms in computer vision, each offering a different trade-off between fidelity, resolution, and computational cost.
The two families most relevant to this study are class-activation mapping, exemplified by Grad-CAM, and game-theoretic feature attribution, exemplified by SHAP, which provide complementary spatial and contribution-based perspectives on a model's behaviour.

\subsubsection{Grad-CAM}
\label{subsubsec:gradcam}

Grad-CAM \citep{selvaraju2017gradcam} produces class-discriminative localisation maps by weighting the activations of a convolutional layer with the gradients of the target class flowing into that layer, yielding a heatmap that highlights the image regions most responsible for the prediction.
Its principal advantages are that it is architecture-agnostic for CNNs, requires no modification or retraining of the model, and is computationally inexpensive, which explains its widespread adoption across application domains.
Successor methods refine this idea: Grad-CAM++ \citep{chattopadhay2018gradcampp} improves localisation when multiple instances of an object are present, while Score-CAM \citep{wang2020scorecam} removes the dependence on gradients altogether by weighting activation maps according to their forward-pass contribution to the class score, often producing cleaner and more stable maps.
In applied work, \citet{kumaran2024lungcam} use Grad-CAM to interpret an ensemble of VGG16, ResNet50, and InceptionV3 for lung-cancer imaging, \citet{ahmed2024clothing} combine multi-backbone transfer learning with Grad-CAM for clothing classification, and \citet{karim2024grapeleaf} pair a lightweight MobileNetV3 edge model with Grad-CAM for grape-leaf diagnosis, collectively demonstrating that the method generalises across architectures and tasks.
A recognised limitation, however, is that Grad-CAM maps are produced at the spatial resolution of a deep convolutional layer and are then upsampled, so they are inherently coarse and can blur fine boundaries, which is a particular concern for fine-grained discrimination.
This coarseness means that Grad-CAM is best treated as an indicator of approximate spatial attention rather than as a precise, pixel-accurate account of the evidence, motivating its combination with a complementary attribution method.

\subsubsection{SHAP}
\label{subsubsec:shap}

SHAP \citep{lundberg2017shap} grounds explanation in cooperative game theory, assigning each feature a Shapley value that quantifies its marginal contribution to a prediction averaged over feature coalitions, which yields a principled measure of positive and negative support.
For images, features are typically defined over regions or superpixels, and SHAP estimates how masking or revealing each region changes the predicted class probability, so the resulting attributions can indicate not only which areas support a class but also which areas argue against it.
This signed, contribution-based view complements the unsigned spatial emphasis of Grad-CAM, since it can reveal suppressive evidence and quantify relative importance rather than merely localising attention.
\citet{nazim2025malware} combine SHAP with LIME and Grad-CAM to interpret malware-image classifiers, illustrating how multiple attribution methods can be triangulated for a more trustworthy explanation.
\citet{tempel2024shapgradcam} directly compare SHAP and Grad-CAM in human-activity recognition and report that the two methods can disagree, with each surfacing aspects of the model's reasoning that the other misses, which reinforces the value of using them jointly.
The principal drawback of SHAP is computational: faithfully approximating Shapley values requires many model evaluations across masked feature coalitions, making image-level SHAP considerably more expensive than a single Grad-CAM pass and dependent on the chosen masking or superpixel scheme.
In this study, SHAP is therefore applied selectively to complement Grad-CAM, providing a region-level account of positive and negative evidence where the additional computational cost is justified.

\subsubsection{Explainability Limitations}
\label{subsubsec:xailimits}

Although attribution maps are visually compelling, a critical reading of the literature warns that an attractive heatmap is not, by itself, evidence that a model reasons correctly.
\citet{geirhos2020shortcut} document how deep networks frequently exploit shortcut features, latching onto incidental correlations that happen to predict the label on the training distribution but fail to generalise, and \citet{ye2024cleverhans} survey the broad prevalence of such spurious correlations across tasks and datasets.
\citet{suhail2025shortcut} show that vision classifiers remain susceptible to this shortcut behaviour, which is especially relevant for sports imagery where backgrounds such as courts, fields, water, or arenas may correlate strongly with the sport and tempt the model to rely on context rather than the athlete or action itself.
\citet{bassi2024lrp} demonstrate, using layer-wise relevance propagation, that explicitly optimising attributions can reduce background bias and steer a network toward foreground evidence, indicating that explanation methods can be diagnostic of, and even corrective for, contextual over-reliance.
A further concern is stability: \citet{kares2025saliency} evaluate saliency maps and find that they can be sensitive to small perturbations and methodological choices, so a single explanation should not be over-interpreted as a definitive account of model behaviour.
\citet{ennab2025pixel} argue for finer, pixel-level interpretability over coarse class-activation maps in high-stakes medical imaging, highlighting that the resolution and faithfulness of an explanation materially affect its trustworthiness.
The consensus emerging from these works is that explanations require careful, domain-informed interpretation, ideally triangulated across complementary methods and corroborated against quantitative robustness evidence, rather than accepted at face value.
This study adopts that cautious stance by pairing Grad-CAM with SHAP and by examining whether highlighted regions correspond to genuine sport-relevant content rather than background shortcuts.

\subsection{Robustness, Generalization, and Deployment}
\label{subsec:robustdeploy}

Beyond clean-data accuracy, a practically useful classifier must remain reliable under the noise, blur, and occlusion that characterise real images, and it must run within realistic resource budgets.
\citet{siedel2023corruption} study robustness under random $L_p$-norm corruptions and show that standard classifiers can degrade markedly when inputs are perturbed, underscoring that high test accuracy on curated data does not guarantee resilience in deployment.
\citet{ramirezagudelo2026noisy} report a complementary finding, that deliberately exposing CNNs to noisy data during training can improve robustness, suggesting that controlled augmentation with realistic corruptions is a viable mitigation strategy rather than a purely harmful nuisance.
Together these results imply that robustness should be evaluated explicitly and, where possible, cultivated through training, rather than assumed from clean-set performance alone.
Deployment imposes a second, orthogonal set of constraints centred on model size, memory, and inference latency, which determine whether a model can run on edge, mobile, or web platforms.
\citet{isong2025lightweight} survey strategies for building efficient lightweight CNNs, and \citet{karim2024grapeleaf} demonstrate a concrete edge deployment in which a compact MobileNetV3 model paired with Grad-CAM runs on constrained hardware, showing that interpretability and efficiency are compatible goals.
\citet{foysal2024plantleaf} similarly couple a CNN classifier with a mobile application, illustrating an end-to-end path from model training to an interactive, user-facing tool of the kind this study targets through a Gradio interface.
These deployment-oriented studies reinforce the rationale for selecting compact backbones such as EfficientNetV2-B0 and ConvNeXt-Tiny, since their favourable accuracy-to-cost profiles make an interpretable, web-deployable sports classifier realistic without sacrificing predictive quality.

\subsection{Comparative Literature Review Matrix}
\label{subsec:litmatrix}

To make the foregoing comparisons explicit, Table~\ref{tab:litmatrix} consolidates the most relevant prior studies along the dimensions that matter most for this work, namely the task domain, the architectures and datasets used, the class granularity, the evaluation metrics reported, whether explainability is incorporated, and whether robustness or deployment is considered.
Reading the matrix across rows rather than down columns clarifies how individual studies cluster, and it exposes the systematic gaps that motivate the present design.
The table juxtaposes sports-specific classification work, general-purpose backbone proposals, applied transfer-learning studies in adjacent domains, and the explainability and robustness literature, so that the contributions and omissions of each can be compared on a common footing.

A first recurring pattern in Table~\ref{tab:litmatrix} is that the sports-classification studies consistently confirm the effectiveness of deep transfer learning and attention, yet they predominantly report top-1 accuracy and rarely include macro-F1, top-$k$ accuracy, or confusion-matrix analysis, which limits insight into fine-grained failure modes.
A second pattern is that explainability and sports image classification are seldom paired: the XAI methods \citep{selvaraju2017gradcam,chattopadhay2018gradcampp,wang2020scorecam,lundberg2017shap} and their applications cluster in medical, security, and agricultural domains \citep{kumaran2024lungcam,nazim2025malware,karim2024grapeleaf,ahmed2024clothing}, while the sports studies \citep{liu2024comparison,li2024seres,zhou2024tuna,fu2024federated,bera2023finegrained} largely omit attribution analysis.
A third pattern is that robustness evaluation, external validation, and efficiency or deployment analysis are scarce within the sports literature, with corruption-robustness and lightweight-deployment evidence \citep{siedel2023corruption,ramirezagudelo2026noisy,isong2025lightweight,foysal2024plantleaf} originating almost entirely outside it.
The proposed study addresses these recurring weaknesses directly: it compares a custom CNN against two complementary pretrained backbones on the same 100-class benchmark, reports a comprehensive metric suite that includes macro-F1 and top-$k$ accuracy alongside confusion-matrix analysis, integrates both Grad-CAM and SHAP to interrogate model reasoning, and explicitly considers robustness and lightweight web deployment.
In doing so, it occupies a position in the matrix that, to the best of our reading, no single prior sports study fully covers, combining comparative architecture evaluation, rigorous metrics, dual-method explainability, and deployment awareness within one coherent pipeline.

\subsection{Synthesis of Literature and Identified Gap}
\label{subsec:synthesis}

Synthesising across the themes above, the field demonstrably does several things well: deep CNNs with transfer learning and attention reliably achieve strong accuracy on sports imagery \citep{liu2024comparison,li2024seres,zhou2024tuna}, modern efficient backbones such as EfficientNetV2 and ConvNeXt deliver high accuracy at a manageable parameter cost \citep{tan2021efficientnetv2,liu2022convnext}, and mature attribution methods exist to probe model decisions \citep{selvaraju2017gradcam,lundberg2017shap}.
The applied transfer-learning literature in adjacent domains further confirms that ImageNet-pretrained features adapt effectively to specialised, data-limited tasks \citep{alotaibi2024gastro,manole2024dermatology}, and the robustness and deployment literature shows that compact, interpretable models can be made resilient and practically deployable \citep{siedel2023corruption,karim2024grapeleaf}.
What remains missing is the integration of these strengths within a single fine-grained sports study: existing sports work seldom couples a controlled, multi-architecture comparison with a comprehensive metric suite, rarely subjects its predictions to complementary explainability analysis, and only infrequently reports robustness or deployment characteristics.
The result is a body of literature that is individually mature along each axis but collectively fragmented, leaving practitioners without a worked example that simultaneously satisfies the demands of accuracy, interpretability, reliability, and usability for many-class sports recognition.
This study is positioned precisely at that intersection, contributing a comparative, metric-rich, dually explainable, and deployment-aware treatment of 100-class sports classification that consolidates the scattered best practices identified above.
Having established this gap and the design principles that follow from it, the next section details the methodology through which these principles are realised.


% ----------- Comparative literature review matrix (Table 1) -----------
\begin{sidewaystable}
\centering
\caption{Comparative literature review matrix summarising datasets, models, evaluation practices,
explainability, robustness, validation, and efficiency analysis across closely related
CNN image-classification studies, contrasted with the proposed study (final row). Column
abbreviations: TL~=~transfer learning, CB~=~class balancing, MC~=~multi-model comparison,
GC~=~Grad-CAM, SH~=~SHAP, RT~=~robustness testing, EV~=~external validation, ST~=~statistical
testing, EA~=~efficiency analysis. Marks: \cmark{}~included, \xmark{}~not included, \pmark{}~partial/unclear.}
\label{tab:litmatrix}
\tiny
\setlength{\tabcolsep}{2.5pt}
\renewcommand{\arraystretch}{1.15}
\begin{tabular}{@{}p{1.5cm}p{1.6cm}p{1.4cm}cp{1.8cm}ccccccccc p{1.5cm}p{2.1cm}@{}}
\toprule
\textbf{Study} & \textbf{Problem} & \textbf{Dataset} & \textbf{Cls} & \textbf{Models} &
\textbf{TL} & \textbf{CB} & \textbf{MC} & \textbf{GC} & \textbf{SH} & \textbf{RT} & \textbf{EV} &
\textbf{ST} & \textbf{EA} & \textbf{Metrics} & \textbf{Key limitation} \\
\midrule
Zhou et al.\ \citep{zhou2024tuna} & Sports image classification & Sports image set & $\sim$18 & Deep net + tuna-swarm opt. & \pmark & \xmark & \cmark & \xmark & \xmark & \xmark & \xmark & \pmark & \pmark & Accuracy & No explainability; opt.-focused \\
Li et al.\ \citep{li2024seres} & Sports image classification & Sports image set & multi & SE-RES-CNN & \pmark & \xmark & \cmark & \xmark & \xmark & \xmark & \xmark & \pmark & \pmark & Acc, Prec, Rec & Accuracy-centric; no XAI \\
Fu et al.\ \citep{fu2024federated} & Sports (federated) & Sports image data & 16 & GA-optimised CNN, federated & \pmark & \xmark & \cmark & \xmark & \xmark & \xmark & \xmark & \pmark & \cmark & Accuracy & Few classes; no XAI \\
Liu \citep{liu2024comparison} & Sports image classification & 100-class sports & 100 & 4 CNNs compared & \cmark & \xmark & \cmark & \xmark & \xmark & \xmark & \xmark & \xmark & \pmark & Accuracy & Accuracy-only; no XAI \\
Bera et al.\ \citep{bera2023finegrained} & Fine-grained sports/yoga/dance & SYD dataset & 102 & SYD-Net (attention) & \cmark & \pmark & \cmark & \pmark & \xmark & \xmark & \pmark & \cmark & Acc, Top-$k$ & No pixel-level XAI \\
Ray et al.\ \citep{ray2025womensports} & Women's sport actions & WSports & multi & ResNet-50 TL & \cmark & \pmark & \cmark & \xmark & \xmark & \xmark & \xmark & \pmark & Acc, F1 & Small data; no XAI \\
Al-Otaibi et al.\ \citep{alotaibi2024gastro} & GI endoscopy & Kvasir & 8 & EfficientNet TL & \cmark & \pmark & \cmark & \xmark & \xmark & \xmark & \pmark & \pmark & Acc, F1, AUC & No XAI; single domain \\
Ahmed et al.\ \citep{ahmed2024clothing} & Clothing/temperature & Clothing images & multi & Multi-backbone TL & \cmark & \pmark & \cmark & \cmark & \xmark & \xmark & \xmark & \pmark & Acc, F1 & No SHAP; no robustness \\
Kumaran et al.\ \citep{kumaran2024lungcam} & Lung cancer & CT images & 3 & Ensemble TL & \cmark & \pmark & \cmark & \cmark & \xmark & \xmark & \xmark & \pmark & Acc, Prec, Rec & Grad-CAM only; no SHAP \\
Nazim et al.\ \citep{nazim2025malware} & Malware imagery & Malware images & multi & CNN + XAI & \cmark & \pmark & \cmark & \cmark & \cmark & \xmark & \xmark & \pmark & Acc, F1 & No robustness/external val. \\
Ram\'irez-A.\ et al.\ \citep{ramirezagudelo2026noisy} & Robustness study & CIFAR-10 & 10 & ResNet-18 & \pmark & \xmark & \cmark & \xmark & \xmark & \cmark & \pmark & \cmark & \pmark & Accuracy & Toy dataset; no XAI \\
Karim et al.\ \citep{karim2024grapeleaf} & Grape leaf disease & Leaf images & 4 & MobileNetV3 (edge) & \cmark & \pmark & \cmark & \cmark & \xmark & \pmark & \xmark & \pmark & \cmark & Acc, F1 & Single crop; no SHAP \\
Foysal et al.\ \citep{foysal2024plantleaf} & Plant leaf disease & Leaf images & multi & CNN + mobile app & \pmark & \pmark & \pmark & \xmark & \xmark & \xmark & \xmark & \xmark & \cmark & Accuracy & No XAI; limited eval. \\
Isong \citep{isong2025lightweight} & Lightweight CNNs & MNIST/Fashion & 10 & Lightweight CNNs & \xmark & \xmark & \cmark & \xmark & \xmark & \xmark & \xmark & \pmark & \cmark & Acc, Size & Simple datasets; no XAI \\
Manole et al.\ \citep{manole2024dermatology} & Dermatology & Skin images & multi & EfficientNetB3 TL & \cmark & \pmark & \pmark & \pmark & \xmark & \xmark & \pmark & \pmark & Acc, AUC & Limited XAI; single domain \\
Bassi et al.\ \citep{bassi2024lrp} & Background bias & X-ray/others & multi & ISNet (LRP) & \cmark & \pmark & \cmark & \cmark & \xmark & \cmark & \cmark & \cmark & Acc, F1, AUC & Complex; not sports \\
\midrule
\textbf{Proposed study} & \textbf{Sports image classification} & \textbf{100 Sports \citep{piosenka2022sports}} & \textbf{100} & \textbf{Custom CNN, EfficientNetV2-B0, ConvNeXt-Tiny} & \cmark & \pmark & \cmark & \cmark & \cmark & \pmark & \xmark & \pmark & \cmark & \textbf{Acc, Top-$k$, macro-F1, efficiency} & Single dataset; no external validation \\
\bottomrule
\end{tabular}
\end{sidewaystable}

\section{Methodology}
\label{sec:methodology}

This section specifies the complete experimental methodology used to build, train, evaluate, and explain a set of convolutional models for fine-grained sports image classification, in sufficient detail to permit independent reproduction.
The study compares a custom convolutional neural network (CNN) baseline against two modern ImageNet-pretrained backbones, EfficientNetV2-B0 \citep{tan2021efficientnetv2} and ConvNeXt-Tiny \citep{liu2022convnext}, on the publicly available ``100 Sports Image Classification'' dataset \citep{piosenka2022sports}.
The methodology is organised so that each design decision is traceable to the implementation: the dataset and its splits, the data-preparation pipeline, the three model architectures, the two-stage transfer-learning training procedure, the evaluation framework, and the explainable artificial intelligence (XAI) analysis.
All preprocessing, augmentation, and model-specific normalisation are embedded inside the model graphs so that the same trained artefact behaves identically during training, validation, and deployment.
A fixed random seed of $42$ is applied throughout to make data ordering, augmentation sampling, and weight initialisation as deterministic as the runtime permits.
Quantitative predictive performance, computational efficiency, and explanation reliability are reported jointly, reflecting the study's aim of identifying not only the most accurate model but also the most trustworthy and deployable one.
The overall pipeline, from raw image folders through preprocessing, model training, evaluation, and explanation, is summarised in Figure~\ref{fig:workflow}.

\subsection{Research Design}
\label{subsec:method-design}

The research design is experimental, quantitative, and comparative, framing sports recognition as a supervised single-label image classification problem over $100$ mutually exclusive sport categories.
Three models are trained and evaluated under an identical data pipeline and an identical evaluation protocol, so that observed differences can be attributed to architecture and training strategy rather than to inconsistent preprocessing or measurement.
The custom CNN serves as a from-scratch baseline that establishes the difficulty of the task without external knowledge, while the two pretrained backbones quantify the benefit of transfer learning from ImageNet.
The design is therefore a controlled comparison in which the dependent variables are predictive accuracy, class-level quality, computational cost, and explanation fidelity, and the independent variable is the model and its training regime.
Each model is treated as a complete, self-contained graph that ingests raw $224\times224\times3$ images and emits a $100$-dimensional class-probability vector, ensuring that the only differences between experiments are the intended ones.
The full experimental workflow, comprising data loading, model construction, two-stage training, multi-metric evaluation, and the Grad-CAM and SHAP explanation stages, is depicted in Figure~\ref{fig:workflow} and is followed identically for every model.
This structure supports a fair head-to-head comparison and yields evidence that is directly interpretable in terms of practical model selection.

\subsection{Dataset Description}
\label{subsec:method-dataset}

\subsubsection{Dataset Source and Accessibility}
\label{subsubsec:method-source}

All experiments use a single, public dataset, the ``100 Sports Image Classification'' collection curated by Piosenka and distributed through Kaggle under the identifier \texttt{gpiosenka/sports-classification} \citep{piosenka2022sports}.
No external or proprietary dataset is introduced, which keeps the study fully reproducible from a single openly licensed source and avoids any dependence on private data.
The dataset is supplied as three directories, \texttt{train}, \texttt{valid}, and \texttt{test}, together with a CSV file that maps each image filename to its class label and split.
This folder-per-class organisation allows the data to be ingested directly with the Keras utility \texttt{image\_dataset\_from\_directory}, which infers integer labels from the directory structure using \texttt{label\_mode='int'}.
Because the source provides an official partitioning, the experimental splits are taken as released rather than resampled, which removes a common source of inconsistency between studies.
The public availability of both the images and the accompanying label CSV means that every result reported here can be regenerated from the original Kaggle release without additional annotation.

\subsubsection{Dataset Composition}
\label{subsubsec:method-composition}

The dataset comprises $13{,}493$ training images, $500$ validation images, and $500$ test images, for a total of $14{,}493$ labelled examples spanning $100$ sport classes.
The validation and test partitions each contain exactly five images per class, giving a clean and perfectly balanced held-out evaluation of five validation and five test images for every one of the $100$ categories.
The training partition contains the remaining images and averages approximately $135$ images per class, which provides a moderate but non-trivial amount of supervision for each category.
Every image is stored as a three-channel JPG at a uniform resolution of $224\times224\times3$ pixels, so no rescaling of spatial dimensions is required to match the input expected by standard ImageNet backbones.
The uniform size and channel count simplify the pipeline considerably, because all images can be batched without padding or aspect-ratio correction.
The combination of a large balanced training set and small balanced evaluation sets is well suited to a controlled comparison, although the modest per-class evaluation count is accounted for explicitly in the statistical-reliability discussion.

\subsubsection{Class Definitions}
\label{subsubsec:method-classes}

The $100$ classes correspond to distinct sport categories, ranging from widely recognised disciplines to visually similar and fine-grained activities that differ only in subtle equipment, posture, or scene cues.
Each image carries exactly one ground-truth label, so the task is single-label multi-class classification rather than multi-label tagging, and the categories are treated as mutually exclusive.
The integer label assigned by \texttt{image\_dataset\_from\_directory} is derived from the alphabetical ordering of the class subdirectories, and this fixed ordering is preserved consistently across training, validation, testing, and all reported confusion matrices.
Fine-grained distinctions among related categories, such as different racquet sports, combat disciplines, or water-based events, make the taxonomy challenging and motivate the use of top-$k$ accuracy alongside top-$1$ accuracy.
The class labels recovered from the directory structure are cross-checked against the provided CSV mapping to confirm that the numeric indices used during training correspond to the intended human-readable sport names.
Maintaining a single canonical class ordering ensures that per-class metrics and confusion-matrix axes are directly comparable across the three models.

\subsubsection{Dataset Quality Assessment}
\label{subsubsec:method-quality}

The dataset is assessed for the properties that most affect supervised image classification, namely consistent dimensionality, valid channel structure, balanced labels, and absence of corrupt files.
All images conform to the stated $224\times224\times3$ JPG format, which removes the variability in size and aspect ratio that often complicates real-world image pipelines.
The held-out partitions are exactly balanced at five images per class, and the training partition is near-balanced at roughly $135$ images per class, so no class is grossly under-represented.
Because the dataset is widely used and curated, gross labelling errors are expected to be rare; nonetheless, the label CSV is used to verify that filenames, directory membership, and class indices remain mutually consistent before training.
Loading the data through the standard Keras utility additionally surfaces any unreadable or malformed image during dataset construction, providing an implicit integrity check on the JPG files.
The principal quality limitation is the small size of the evaluation partitions, which constrains the precision of per-class estimates and is therefore treated as a known caveat rather than a defect requiring correction.

\subsubsection{Dataset Visualization}
\label{subsubsec:method-visualization}

To document the visual character of the data and to confirm correct label alignment, a representative grid of training images with their associated class labels is generated and shown in Figure~\ref{fig:samples}.
This sample grid illustrates the wide variation in viewpoint, scale, lighting, and background that the models must accommodate, as well as the visual similarity between certain fine-grained categories.
The empirical distribution of training images across the $100$ classes is plotted in Figure~\ref{fig:dist}, which makes the near-balanced nature of the training set directly visible.
The class-distribution plot confirms that per-class counts cluster tightly around the mean of approximately $135$ images, with no class exhibiting the extreme scarcity that would necessitate resampling.
Together these two visualisations verify that the data have been loaded with correct labels and that the balance assumption underpinning the class-imbalance strategy is justified by the actual counts.
Inspecting the data visually before modelling also guards against silent pipeline errors, such as mismatched labels or unintended channel reordering, that summary statistics alone might miss.

\subsection{Data Preparation}
\label{subsec:method-prep}

\subsubsection{Data Cleaning}
\label{subsubsec:method-cleaning}

Because the dataset is pre-curated and uniformly formatted, data cleaning is light and is restricted to verification rather than transformation.
Images are confirmed to share the expected $224\times224\times3$ shape, and any file that fails to decode during dataset construction would be excluded by the loader, ensuring that only valid JPG images enter the pipeline.
The label CSV is used to confirm that every image filename maps to a single, valid class and split, so that ambiguous or duplicated label assignments are detected before training.
No manual relabelling, cropping, or deduplication is applied, since the official curation already enforces consistent structure and the study deliberately preserves the released data to maximise reproducibility.
Pixel values are left in their native $[0,255]$ range at this stage, deferring all normalisation to the model-specific preprocessing layers described below.
This minimal-intervention cleaning policy keeps the input identical across all three models and avoids introducing preprocessing artefacts that could confound the comparison.

\subsubsection{Data Splitting}
\label{subsubsec:method-split}

The study uses the official train/validation/test partitioning supplied with the dataset, comprising $13{,}493$ training, $500$ validation, and $500$ test images, and does not re-split or shuffle examples across partitions.
Adopting the released splits guarantees that no image appears in more than one partition, eliminating data leakage between training and evaluation by construction.
The validation partition is used exclusively for monitoring during training, for early stopping, and for learning-rate scheduling, while the test partition is reserved strictly for final, single-pass evaluation of each trained model.
Keeping the test set untouched until the final evaluation ensures that reported test metrics reflect genuine generalisation rather than information that leaked through repeated tuning.
Because the partitions are class-balanced in their held-out portions, the validation and test estimates are not biased toward any particular sport category.
This strict separation of training, validation, and test roles is applied identically to all three models so that their reported test results are mutually comparable.

\subsubsection{Image Preprocessing}
\label{subsubsec:method-preprocessing}

Images are loaded and retained in the raw integer range $[0,255]$, and each model applies its own model-appropriate normalisation internally rather than relying on a shared external transform.
The custom CNN includes a \texttt{Rescaling} layer that divides pixel values by $255$ to map inputs into $[0,1]$ as its first numeric operation after augmentation.
The two pretrained backbones instead apply their official preprocessing functions inside the graph, namely \texttt{efficientnet\_v2.preprocess\_input} for EfficientNetV2-B0 and \texttt{convnext.preprocess\_input} for ConvNeXt-Tiny, so that inputs match the statistics expected by the ImageNet-pretrained weights.
Embedding preprocessing as layers within each model means that the exported artefact normalises inputs consistently, removing any risk of a train-test preprocessing mismatch.
The target spatial size is fixed at $224\times224$, which coincides with both the native resolution of the dataset and the canonical input size of the chosen backbones, so no resizing is necessary.
All images are processed as three-channel RGB tensors, and the batched dataset feeds these tensors directly into the model-specific preprocessing stage.

\subsubsection{Data Augmentation}
\label{subsubsec:method-augmentation}

Data augmentation is implemented as a sequence of Keras preprocessing layers placed at the front of every model and active only at training time, automatically disabled during validation, testing, and inference.
The augmentation pipeline applies \texttt{RandomFlip('horizontal')}, \texttt{RandomRotation(0.08)}, \texttt{RandomZoom(0.10)}, and \texttt{RandomContrast(0.10)}, introducing mild geometric and photometric variation that discourages overfitting to incidental scene details.
Only horizontal flipping is used, because the identity of a sport is invariant to mirroring, whereas vertical flipping is deliberately excluded as it would produce physically implausible, upside-down scenes that do not occur in real sports imagery.
The rotation, zoom, and contrast ranges are kept small so that augmented images remain realistic and the underlying sport remains recognisable, providing regularisation without distorting the class signal.
Because the augmentation layers live inside the model graph, the same augmentation behaviour is applied identically to all three architectures, keeping the regularisation contribution consistent across the comparison.
Performing augmentation on-graph also ensures that each training epoch samples fresh transformations of every image, effectively enlarging the apparent diversity of the moderate-sized training set.

\subsubsection{Class-Imbalance Handling}
\label{subsubsec:method-imbalance}

The training partition is near-balanced, with approximately $135$ images per class and no category exhibiting severe under-representation, as confirmed by the distribution shown in Figure~\ref{fig:dist}.
Given this near-uniform class frequency, no resampling, class weighting, or oversampling is applied, because such corrections offer little benefit when the classes are already comparable in size and can introduce their own distortions.
The mild augmentation described above provides modest, class-agnostic regularisation that benefits all categories uniformly rather than artificially inflating any minority class.
Because the held-out validation and test partitions are exactly balanced at five images per class, the macro-averaged metrics reported later are not biased by differential class frequency at evaluation time.
Avoiding resampling also keeps the effective training distribution faithful to the released dataset, which strengthens the reproducibility of the reported results.
This decision is revisited in the evaluation framework, where macro-averaging is used precisely so that every class contributes equally to the headline quality metrics regardless of small residual count differences.

\subsection{Model Development}
\label{subsec:method-models}

\subsubsection{Custom CNN Baseline}
\label{subsubsec:method-customcnn}

The custom CNN is a compact from-scratch baseline that begins with an \texttt{Input(224,224,3)} layer, followed by the shared augmentation block and a \texttt{Rescaling(1/255)} layer that maps pixels into $[0,1]$.
The feature extractor consists of four convolutional blocks with progressively increasing filter counts of $32$, $64$, $128$, and $256$, where each block contains a $3\times3$ \texttt{Conv2D} layer with same padding and ReLU activation, followed by \texttt{BatchNormalization} and \texttt{MaxPooling2D}.
This progression deepens the channel dimension while halving the spatial resolution at each stage, yielding an increasingly abstract representation in the manner characteristic of conventional CNNs.
After the final block, a \texttt{GlobalAveragePooling2D} layer collapses the spatial dimensions into a single feature vector, reducing parameters and providing a degree of translation invariance.
A classification head then applies \texttt{Dropout(0.4)}, a \texttt{Dense(256)} layer with ReLU activation, a further \texttt{Dropout(0.3)}, and finally a \texttt{Dense(100)} layer with softmax activation that outputs the $100$-class probability distribution.
The two dropout layers and the batch-normalisation within each block provide regularisation that is important when training a network of this capacity without pretrained weights.
This baseline establishes the level of performance attainable using only the dataset itself, serving as the reference against which the transfer-learning models are judged.

\subsubsection{Transfer Learning: EfficientNetV2-B0}
\label{subsubsec:method-effnet}

The first transfer-learning model is built on EfficientNetV2-B0 \citep{tan2021efficientnetv2}, a compound-scaled convolutional backbone designed for a strong accuracy-to-efficiency trade-off.
The model graph starts with an \texttt{Input} layer and the shared augmentation block, after which \texttt{efficientnet\_v2.preprocess\_input} normalises the inputs to the distribution expected by the pretrained weights.
The backbone is instantiated as \texttt{EfficientNetV2B0(include\_top=False, weights='imagenet')}, so that its ImageNet-pretrained convolutional features are reused while its original classification head is discarded.
A lightweight classification head is appended, consisting of \texttt{GlobalAveragePooling2D}, \texttt{Dropout(0.3)}, and a \texttt{Dense(100)} layer with softmax activation that maps the pooled features onto the $100$ sport classes.
Keeping the head deliberately small concentrates the model's adaptation to the target domain in a few trainable parameters during the initial feature-extraction phase, reducing the risk of overfitting.
This configuration leverages the broad visual knowledge captured by ImageNet pretraining while requiring only modest additional training on the sports data.

\subsubsection{Transfer Learning: ConvNeXt-Tiny}
\label{subsubsec:method-convnext}

The second transfer-learning model uses ConvNeXt-Tiny \citep{liu2022convnext}, a modernised convolutional architecture that incorporates design elements inspired by vision transformers while remaining fully convolutional.
Its graph mirrors the EfficientNetV2 model exactly except for the backbone and its matching normalisation: inputs pass through the shared augmentation block and then \texttt{convnext.preprocess\_input} before entering the backbone.
The backbone is instantiated as \texttt{ConvNeXtTiny(include\_top=False, weights='imagenet')}, again reusing ImageNet-pretrained features and discarding the original classifier.
The identical head, comprising \texttt{GlobalAveragePooling2D}, \texttt{Dropout(0.3)}, and a softmax \texttt{Dense(100)} layer, is attached so that the only methodological difference from the EfficientNetV2 model is the backbone itself.
Using a shared head and a shared augmentation pipeline across both pretrained models ensures that the comparison between EfficientNetV2-B0 and ConvNeXt-Tiny isolates the effect of the backbone architecture.
This design allows the study to assess whether the transformer-inspired ConvNeXt design offers any advantage over compound-scaled EfficientNetV2 on the fine-grained sports task.

\subsubsection{Feature Extraction and Fine-Tuning Strategy}
\label{subsubsec:method-finetune}

The two pretrained models are trained with a two-stage transfer-learning strategy that first extracts features and then selectively fine-tunes the backbone.
In Stage~1, the entire backbone is frozen and only the newly added classification head is trained, allowing the randomly initialised head to adapt to the pretrained features at a relatively high learning rate of $10^{-3}$ for up to eight epochs.
In Stage~2, the top approximately $30\%$ of the backbone layers are unfrozen while the bottom $70\%$ remain frozen, and the model is fine-tuned at a much smaller learning rate of $10^{-5}$ for up to eight further epochs.
Unfreezing only the upper portion of the backbone adapts the most task-specific, high-level features to the sports domain while preserving the general-purpose low-level features that ImageNet pretraining provides.
The small Stage~2 learning rate is essential to avoid catastrophically overwriting the pretrained weights, since large updates at this point would erase the very knowledge that transfer learning is intended to exploit.
The custom CNN does not undergo this two-stage procedure; instead it is trained from random initialisation in a single stage of up to $30$ epochs at a learning rate of $10^{-3}$, reflecting its lack of any pretrained weights to preserve.
This staged strategy balances rapid head adaptation against careful, conservative refinement of the transferred representation.

\subsection{Training Procedure}
\label{subsec:method-training}

\subsubsection{Loss Function}
\label{subsubsec:method-loss}

All models are optimised with the sparse categorical cross-entropy loss, which is appropriate for single-label multi-class classification with integer-encoded targets and a softmax output.
For a single example with integer ground-truth class $y$ and predicted class-probability vector $\hat{p}$ over $C=100$ classes, the loss is defined in Equation~\eqref{eq:loss}.
\begin{equation}
\label{eq:loss}
\mathcal{L} = -\log\left(\hat{p}_{y}\right),
\end{equation}
where $\hat{p}_{y}$ is the predicted probability assigned to the true class, and the reported training objective is the mean of this quantity over the mini-batch.
The sparse formulation consumes the integer labels produced by \texttt{label\_mode='int'} directly, avoiding the need to one-hot encode the $100$ classes and thereby reducing memory overhead.
Minimising this loss is equivalent to maximising the log-likelihood of the correct sport class under the model, which aligns the training objective with the top-$1$ accuracy of primary interest.
The same loss is used unchanged across the custom CNN and both transfer-learning models so that optimisation targets are identical in every experiment.

\subsubsection{Optimizer and Learning Rate}
\label{subsubsec:method-optimizer}

All models are trained with the Adam optimiser, whose adaptive per-parameter learning rates provide stable and efficient convergence across the differing scales of the three architectures.
The base learning rate is set to $10^{-3}$ for the custom CNN and for Stage~1 of both transfer-learning models, which is a standard, comparatively aggressive rate suitable for training a freshly initialised classification head.
For Stage~2 fine-tuning of the pretrained backbones, the learning rate is lowered to $10^{-5}$ so that the unfrozen high-level layers are adjusted gently and the transferred representation is preserved.
This two-tier learning-rate scheme is central to the transfer-learning strategy, separating the fast adaptation of the new head from the slow, careful refinement of the pretrained features.
The learning rate is further modulated during training by an on-plateau scheduler, described below, which reduces it automatically when validation performance stops improving.
Using a single optimiser family with these clearly defined learning rates keeps the optimisation procedure transparent and reproducible.

\subsubsection{Training Controls}
\label{subsubsec:method-controls}

Training is governed by a consistent set of controls applied to every model, comprising a batch size of $32$, an image size of $224$, and the global random seed of $42$.
An \texttt{EarlyStopping} callback monitors validation accuracy with a patience of four epochs and restores the best weights observed, so that training halts once validation accuracy ceases to improve and the model is rolled back to its peak state.
A \texttt{ReduceLROnPlateau} callback monitors validation loss with a factor of $0.5$, a patience of two epochs, and a floor of \texttt{min\_lr}~$=10^{-6}$, halving the learning rate whenever the validation loss stalls.
Together these callbacks prevent both overfitting from excessive training and stagnation from an inappropriately high learning rate, while the epoch caps of eight per stage for transfer learning and $30$ for the custom CNN bound the total training budget.
Restoring the best weights ensures that the evaluated model corresponds to the epoch of strongest validation performance rather than the final, potentially overfit epoch.
Because these controls are identical across all models, differences in their training behaviour reflect architecture and initialisation rather than inconsistent stopping or scheduling.

\subsubsection{Hyperparameter Selection}
\label{subsubsec:method-hyper}

The hyperparameters were chosen to be standard, defensible defaults for transfer learning and compact CNNs rather than the product of an exhaustive search, in keeping with the study's emphasis on a fair and reproducible comparison.
The complete set of training hyperparameters, including optimiser, loss, batch size, learning rates, epoch budgets, augmentation ranges, and callback settings, is consolidated in Table~\ref{tab:hyper}.
The same values are applied to all three models wherever they are shared, with the only intentional differences being the per-model preprocessing, the backbone, and the two-stage versus single-stage training regime.
This disciplined, mostly shared hyperparameter configuration ensures that the comparative results are attributable to the models themselves rather than to differential tuning effort.

\begin{table}[htbp]
\centering
\caption{Training and pipeline hyperparameters used for all models.}
\label{tab:hyper}
\begin{tabular}{ll}
\toprule
\textbf{Hyperparameter} & \textbf{Value} \\
\midrule
Input image size & $224\times224\times3$ \\
Batch size & $32$ \\
Random seed & $42$ \\
Optimiser & Adam \\
Loss function & Sparse categorical cross-entropy \\
Custom CNN learning rate & $10^{-3}$ \\
Custom CNN max epochs & $30$ \\
Transfer Stage~1 learning rate & $10^{-3}$ \\
Transfer Stage~1 max epochs & $8$ (backbone frozen) \\
Transfer Stage~2 learning rate & $10^{-5}$ \\
Transfer Stage~2 max epochs & $8$ (top ${\sim}30\%$ unfrozen) \\
EarlyStopping monitor & Validation accuracy \\
EarlyStopping patience & $4$ (restore best weights) \\
ReduceLROnPlateau monitor & Validation loss \\
ReduceLROnPlateau factor & $0.5$ \\
ReduceLROnPlateau patience & $2$ \\
Minimum learning rate & $10^{-6}$ \\
RandomFlip & Horizontal \\
RandomRotation & $0.08$ \\
RandomZoom & $0.10$ \\
RandomContrast & $0.10$ \\
\bottomrule
\end{tabular}
\end{table}

\subsection{Evaluation Framework}
\label{subsec:method-eval}

\subsubsection{Primary Predictive Metrics}
\label{subsubsec:method-primary}

The primary predictive metrics are computed on the held-out test set and quantify how accurately each model assigns the correct sport category among the $C=100$ classes.
Overall accuracy is the fraction of correctly classified test images, defined in Equation~\eqref{eq:acc}, where $N$ is the number of test images, $y_i$ is the true class of image $i$, $\hat{y}_i$ is the predicted class, and $\mathbb{1}[\cdot]$ is the indicator function.
\begin{equation}
\label{eq:acc}
\mathrm{Accuracy} = \frac{1}{N}\sum_{i=1}^{N}\mathbb{1}\!\left[\hat{y}_i = y_i\right].
\end{equation}
For each class $c$, precision and recall are computed from the per-class true positives $\mathrm{TP}_c$, false positives $\mathrm{FP}_c$, and false negatives $\mathrm{FN}_c$, and the macro-averaged precision and recall over the $100$ sports are given by Equation~\eqref{eq:macropr}.
\begin{equation}
\label{eq:macropr}
\mathrm{Precision}_{\mathrm{macro}} = \frac{1}{C}\sum_{c=1}^{C}\frac{\mathrm{TP}_c}{\mathrm{TP}_c+\mathrm{FP}_c}, \qquad
\mathrm{Recall}_{\mathrm{macro}} = \frac{1}{C}\sum_{c=1}^{C}\frac{\mathrm{TP}_c}{\mathrm{TP}_c+\mathrm{FN}_c}.
\end{equation}
The macro F1 score is the unweighted mean of the per-class F1 scores, where each class F1 is the harmonic mean of its precision and recall, as expressed in Equation~\eqref{eq:macrof1}.
\begin{equation}
\label{eq:macrof1}
\mathrm{F1}_{\mathrm{macro}} = \frac{1}{C}\sum_{c=1}^{C}\frac{2\,\mathrm{Precision}_c\,\mathrm{Recall}_c}{\mathrm{Precision}_c+\mathrm{Recall}_c}.
\end{equation}
Because the $100$ sport classes are fine-grained and frequently confusable, top-$k$ accuracy is also reported, counting a prediction as correct if the true class lies among the $k$ highest-scoring classes, as defined in Equation~\eqref{eq:topk} where $\mathrm{rank}_i(y_i)$ is the rank of the true class for image $i$ under the model's sorted output scores.
\begin{equation}
\label{eq:topk}
\mathrm{Top}\text{-}k\;\mathrm{Accuracy} = \frac{1}{N}\sum_{i=1}^{N}\mathbb{1}\!\left[\mathrm{rank}_i(y_i) \le k\right].
\end{equation}
Top-$3$ and top-$5$ accuracy are reported specifically because, in a $100$-class sports taxonomy, a model that ranks the correct discipline among its top few guesses is still operationally useful for assisted tagging even when its single top guess is wrong.

\subsubsection{Class-Level Evaluation}
\label{subsubsec:method-classlevel}

Beyond the aggregate metrics, the framework reports class-level quality so that strengths and weaknesses across the $100$ sports can be examined individually.
A per-class F1 score is computed for every category, exposing which sports are reliably recognised and which are systematically confused, information that a single accuracy figure necessarily hides.
A weighted F1 score, which averages per-class F1 values in proportion to the number of true instances of each class, is reported alongside the macro F1 to summarise quality while accounting for any residual differences in class support.
A full $100\times100$ confusion matrix is generated for each model, with rows indexing true classes and columns indexing predicted classes under the fixed canonical class ordering, making clusters of mutual confusion between visually similar sports directly visible.
Macro-averaging is emphasised throughout because it treats every sport equally and therefore reflects performance on rare or difficult categories rather than being dominated by easy, well-populated ones.
This class-level analysis turns the evaluation from a single ranking of models into a diagnostic account of where each model succeeds and where it fails.

\subsubsection{Computational Efficiency}
\label{subsubsec:method-efficiency}

The framework measures computational efficiency so that predictive quality can be weighed against the practical cost of deployment.
For each model the total number of parameters is recorded, and the on-disk model size in megabytes is estimated as the parameter count multiplied by four bytes per parameter and divided by $10^{6}$, as expressed in Equation~\eqref{eq:size}.
\begin{equation}
\label{eq:size}
\mathrm{Size}_{\mathrm{MB}} = \frac{\mathrm{Parameters}\times 4}{10^{6}}.
\end{equation}
Wall-clock training time is logged for each model so that the cost of the from-scratch baseline and the two-stage transfer-learning runs can be compared.
Inference latency is reported as milliseconds per image, measured as the forward-pass time over a single batch divided by the batch size, providing a representative estimate of per-image throughput at deployment.
Reporting parameters, model size, training time, and inference latency together gives a multi-dimensional view of efficiency that complements the accuracy-oriented metrics.
These measurements are essential for the study's deployability assessment, since a marginally more accurate model may be undesirable if it is substantially larger or slower at inference.

\subsubsection{Statistical Reliability}
\label{subsubsec:method-reliability}

The statistical reliability of the reported metrics is interpreted in light of the small evaluation partitions, which contain only five test images per class and $500$ images in total.
With so few examples per class, individual per-class F1 scores are coarse and sensitive to single misclassifications, so they are read as indicative rather than as precise point estimates.
The macro and weighted aggregates over all $100$ classes are correspondingly more stable than any single per-class figure, because they pool information across the full $500$-image test set.
To improve determinism and comparability, all experiments use the fixed seed of $42$, which controls data ordering, augmentation sampling, and weight initialisation as far as the runtime allows.
The confusion matrix is used as a qualitative reliability check, since consistent, structured confusion between specific sport pairs is more trustworthy evidence than isolated errors that may reflect sampling noise.
These considerations frame the headline numbers as well-supported comparative indicators while transparently acknowledging the precision limits imposed by the small held-out sets.

\subsection{Explainable AI Methodology}
\label{subsec:method-xai}

\subsubsection{Grad-CAM Procedure}
\label{subsubsec:method-gradcam}

Gradient-weighted Class Activation Mapping (Grad-CAM) \citep{selvaraju2017gradcam} is applied to the best-performing model to visualise the spatial evidence underlying its predictions.
The method operates on the last four-dimensional convolutional feature map of the model, computing the gradient of the predicted-class score with respect to that feature map to obtain channel-wise importance weights.
These weights are used to form a weighted combination of the feature-map channels, after which a ReLU is applied to retain only the regions that positively support the predicted class.
The resulting coarse heatmap is normalised, upsampled to the $224\times224$ input resolution, and overlaid on the original image so that the model's region of attention is directly visible.
Grad-CAM maps are generated for three diagnostic cases, namely correctly classified high-confidence examples, correctly classified low-confidence examples, and misclassified examples, so that attention can be inspected across the full spectrum of model behaviour.
This stratified selection reveals whether correct predictions are driven by the athlete and equipment rather than by spurious background cues, and whether errors coincide with attention drifting to irrelevant regions.

\subsubsection{SHAP Procedure}
\label{subsubsec:method-shap}

SHapley Additive exPlanations (SHAP) \citep{lundberg2017shap} provide a complementary, attribution-based explanation that assigns a contribution to each input region for the predicted class.
SHAP is computed with the \texttt{GradientExplainer}, using a random background set of $25$ images to approximate the expected model output against which contributions are measured.
The explainer is configured to explain only the top-$1$ predicted class, with \texttt{ranked\_outputs=1} and \texttt{nsamples=50}, which keeps the computation tractable while focusing the attribution on the class the model actually predicts.
Because gradient-based SHAP can be sensitive to the specific runtime and library versions, an occlusion-sensitivity fallback is provided that systematically slides a $32$-pixel grey patch across the image with a stride of $24$ pixels and records the drop in the predicted-class score at each position.
This occlusion procedure produces a model-agnostic importance map that requires only forward passes, ensuring that an explanation is always available even if the gradient-based SHAP computation fails on the execution environment.
Together, the SHAP attributions and the occlusion fallback corroborate or challenge the spatial evidence highlighted by Grad-CAM, strengthening confidence in the interpretation of model behaviour.

\subsubsection{Explanation Assessment}
\label{subsubsec:method-explanation-assessment}

The two explanation methods are assessed jointly and qualitatively, by examining whether Grad-CAM and SHAP attribute predictions to plausible, sport-relevant evidence.
Agreement between the gradient-based attention of Grad-CAM and the attribution maps of SHAP or occlusion increases confidence that a prediction rests on genuine visual cues rather than on dataset artefacts or background shortcuts.
The deliberate inclusion of low-confidence and misclassified cases allows the analysis to characterise failure modes, for example attention that fixates on crowds, logos, or backgrounds instead of on the athlete or equipment.
This assessment directly serves the study's trustworthiness objective, since a model that is both accurate and demonstrably attentive to meaningful evidence is preferable to one that is merely accurate.
Using two methodologically distinct explanation techniques guards against the idiosyncrasies of any single method, providing a more robust basis for judging interpretability.
The explanation analysis thereby connects the quantitative metrics to a qualitative understanding of why each model behaves as it does.

\subsection{Robustness Analysis}
\label{subsec:method-robustness}

A robustness analysis is included as a planned, optional component to probe how stable the models remain when test images are degraded in ways that mimic real-world capture conditions.
The procedure perturbs the test images with controlled corruptions, namely additive Gaussian noise, blur, brightness changes, and partial occlusion, and then re-measures classification accuracy on the perturbed inputs.
Comparing accuracy on clean and corrupted test sets quantifies each model's sensitivity to common image degradations and indicates which architecture degrades most gracefully.
This analysis is relevant to deployment because operational sports imagery is frequently affected by motion blur, variable lighting, compression artefacts, and partial occlusion of the athlete.
Because it is treated as an optional extension, the robustness study augments rather than replaces the core clean-test evaluation, and its results are interpreted as supplementary evidence of practical reliability.
A model that maintains accuracy under such corruptions is more credible for real-world tagging than one whose performance collapses under modest perturbation.

\subsection{Software, Hardware, and Reproducibility}
\label{subsec:method-software}

The implementation is written in Python and built primarily on TensorFlow and Keras for model construction and training, with scikit-learn used for the computation of precision, recall, F1, and confusion matrices.
Explanations are produced with the SHAP library and visualised, together with all figures, using Matplotlib, while the dataset is ingested through the Keras \texttt{image\_dataset\_from\_directory} utility.
All experiments are executed on Kaggle GPU instances, specifically NVIDIA T4 or P100 accelerators, which provide sufficient memory and throughput for training the compact backbones at a batch size of $32$.
Reproducibility is supported by fixing the global random seed to $42$, by relying on the dataset's official splits, and by embedding all preprocessing and augmentation inside the model graphs so that behaviour is identical across runs and environments.
Trained models, a consolidated \texttt{results.json} file containing the recorded metrics, and all generated figures are written to the \texttt{/kaggle/working} directory, and the complete code is released through a public GitHub repository.
This combination of a single public dataset, fixed seeds, open tooling, and published artefacts allows the entire study to be reproduced end to end.

\subsection{Ethical and Practical Considerations}
\label{subsec:method-ethics}

The study uses only publicly available, appropriately licensed images and processes no personal data, so the privacy footprint of the work is minimal.
Nonetheless, the imagery may carry demographic and scene biases inherited from how sports are photographed and which events are represented, and such biases could in principle skew recognition toward better-represented contexts.
Sports tagging is a comparatively low-risk application, since errors lead to mislabelled or misretrieved media rather than to consequential decisions about individuals, which limits the potential for direct harm.
Even so, human oversight is advised when such a system is deployed, so that automated tags are reviewed before they influence high-visibility archives, recommendations, or editorial workflows.
The integration of Grad-CAM and SHAP explanations further supports responsible use by allowing practitioners to inspect whether predictions rest on meaningful visual evidence rather than on spurious correlations.
By foregrounding transparency, bias awareness, and human-in-the-loop oversight, the methodology aligns the technical contribution with practical and ethical expectations for deployed visual classifiers.


% ---- methodology / dataset figures ----
\projfig{figures/workflow.png}{0.95}{End-to-end workflow of the proposed explainable sports-image
classification framework: dataset acquisition and inspection, preprocessing and augmentation, custom
CNN and transfer-learning model development, two-stage training, multi-metric evaluation, Grad-CAM
and SHAP interpretation, and web deployment.}{fig:workflow}

\projfig{figures/sample_grid.png}{0.95}{Representative training images, one per class for a random
subset of the 100 sports categories, illustrating the visual diversity and the fine-grained
similarity between some sports.}{fig:samples}

\projfig{figures/class_distribution.png}{0.9}{Per-class training-image counts across the 100 sports
categories (sorted), showing that the dataset is approximately balanced with roughly 135 images per
class.}{fig:dist}

\section{Results}
\label{sec:results}

This section reports the empirical findings obtained by executing the experimental pipeline described in Section~\ref{sec:methodology} on the ``100 Sports Image Classification'' dataset, covering dataset preparation, training convergence, overall and class-level predictive performance, explainability, robustness, and deployment efficiency.
Numeric values shown in red are placeholders to be replaced by the executed notebook's \texttt{results.json}.
Throughout, evidence is presented with minimal interpretation, and the explanation of \emph{why} the observed patterns arise is deferred to the Discussion in Section~\ref{sec:discussion}.
Results are organised to address the five research questions in turn, beginning with the data foundation that underpins every subsequent measurement and ending with a consolidated table that maps each research question to its principal finding.
All three models, the custom CNN baseline, EfficientNetV2-B0, and ConvNeXt-Tiny, are evaluated under the identical preprocessing and evaluation protocol so that the reported differences can be attributed to architecture and training strategy rather than measurement inconsistency.
Predictive quality is summarised with top-$1$ accuracy and macro-averaged $F_1$, supplemented by top-$3$ and top-$5$ accuracy, while computational cost is summarised with parameter count, on-disk model size, and per-image inference latency.
Explanation behaviour is assessed qualitatively through Grad-CAM and SHAP visualisations and through their agreement, and deployment feasibility is demonstrated with a working interactive prototype.

\subsection{Dataset and Preprocessing Results}
\label{subsec:res-dataset}

After ingestion through the Keras pipeline, the dataset yielded $13{,}493$ usable training images, $500$ validation images, and $500$ test images, for a total of $14{,}493$ labelled examples across the $100$ sport classes, with no corrupt or off-specification files discarded during the integrity check.
Every retained image conformed to the expected $224\times224\times3$ three-channel JPG format, so the final usable count matched the released dataset exactly and no spatial rescaling or channel correction was required before batching.
The class-distribution analysis, plotted in Figure~\ref{fig:dist}, confirmed that the training partition is near-balanced, averaging approximately $135$ images per class, with the most and least populated categories differing only modestly and no class falling into severe under-representation.
The validation and test partitions are exactly balanced, containing precisely five images per class each, which gives a clean and uniform basis for held-out evaluation even though the small per-class evaluation count limits the resolution of class-level estimates.
The representative sample grid in Figure~\ref{fig:samples} verifies correct label alignment and illustrates the visual heterogeneity of the data, including wide variation in viewpoint, scale, lighting, crowd background, and the presence of multiple athletes or pieces of equipment within a single frame.
The same grid also makes the fine-grained nature of the task apparent, since several distinct classes share near-identical scenes, equipment, and postures that differ only in subtle cues.
The on-graph augmentation pipeline, comprising horizontal flip, mild rotation, zoom, and contrast jitter, was confirmed to activate only at training time and to leave validation, test, and inference images untouched, preserving the integrity of the held-out measurements.
Taken together, these preprocessing results establish a balanced, high-integrity data foundation on which the model comparison underlying RQ1 can be conducted without confounding from class imbalance or inconsistent input formatting.

\subsection{Training and Convergence Results}
\label{subsec:res-training}

The training and validation learning curves for all three models are presented in Figure~\ref{fig:curves}, which plots loss and accuracy against epoch for the custom CNN single-stage run and for the two-stage transfer-learning runs of the pretrained backbones.
The custom CNN converged slowly from random initialisation and was halted by early stopping after approximately \res{22} epochs of its $30$-epoch budget, reaching a plateau in validation accuracy while its training accuracy continued to rise, indicating a mild but visible degree of overfitting.
The transfer-learning models converged substantially faster: during Stage~1, with the backbone frozen and only the classification head trained at a learning rate of $10^{-3}$, validation accuracy rose steeply within the first few epochs as the randomly initialised head adapted to the pretrained features.
The transition to Stage~2 fine-tuning, in which the top $30\%$ of backbone layers were unfrozen at a learning rate of $10^{-5}$, produced a further smaller improvement in validation accuracy without destabilising the loss, confirming that the conservative fine-tuning rate preserved the transferred representation.
EfficientNetV2-B0 reached its early-stopping point after roughly \res{12} effective epochs across both stages, and ConvNeXt-Tiny after roughly \res{13}, both well inside their combined sixteen-epoch budget.
The on-plateau learning-rate scheduler was observed to reduce the rate once validation performance stalled, after which the curves flattened with negligible further change, and the best-epoch weights were restored at termination.
The gap between training and validation curves was narrow for the pretrained models and wider for the custom CNN, consistent with stronger regularisation from the transferred features and the limited capacity of the from-scratch baseline to generalise from a moderate training set.
No evidence of catastrophic divergence, exploding loss, or severe underfitting was observed for any model, and all three runs terminated cleanly via early stopping with their best validation weights retained for final test-set evaluation.

\subsection{Overall Model Performance}
\label{subsec:res-overall}

The headline test-set performance of the three models, evaluated once on the held-out $500$-image test partition, is consolidated in Table~\ref{tab:results}, and the corresponding multi-metric comparison is visualised in Figure~\ref{fig:cmp}.
On top-$1$ accuracy, \bestModel{} achieved the strongest result at \res{91.0}\%, ahead of ConvNeXt-Tiny at \res{89.1}\% and the custom CNN baseline at \res{81.3}\%, and the same ordering held for macro-averaged $F_1$ at \res{89.5}\%, \res{88.0}\%, and \res{79.6}\% respectively.
The macro-$F_1$ values track the accuracy values closely for every model, which is expected given the near-balanced evaluation partition, and indicates that the reported accuracies are not inflated by a small number of dominant classes.
The top-$3$ and top-$5$ accuracies were uniformly high, reaching \res{97.6}\% and \res{99.0}\% for \bestModel{}, \res{96.5}\% and \res{98.4}\% for ConvNeXt-Tiny, and \res{92.8}\% and \res{96.0}\% for the custom CNN, showing that the correct class is almost always within the top few predictions even when it is not ranked first.
The two ImageNet-pretrained backbones outperformed the from-scratch custom CNN by a clear margin of roughly \res{8}--\res{10} accuracy points, which is the central quantitative evidence bearing on RQ2 regarding the benefit of transfer learning over a custom architecture.
Figure~\ref{fig:cmp} renders this ordering across all reported metrics simultaneously and shows that the advantage of the pretrained models is consistent rather than confined to a single metric.
The margin between the two pretrained backbones themselves is comparatively small, so the identity of the single best model is determined as much by efficiency as by accuracy, a trade-off examined separately in Section~\ref{subsec:res-efficiency}.
On predictive performance alone, \bestModel{} is identified as the strongest model in this study.

\begin{table}[htbp]
\centering
\caption{Overall test-set performance and computational cost of the three models on the $500$-image held-out test partition. Accuracy, macro-$F_1$, top-$3$, and top-$5$ are percentages; inference latency is per-image. Values in red are placeholders pending execution of the notebook.}
\label{tab:results}
\resizebox{\textwidth}{!}{%
\begin{tabular}{lccccccc}
\toprule
Model & Accuracy (\%) & Macro-$F_1$ (\%) & Top-3 (\%) & Top-5 (\%) & Params (M) & Size (MB) & Inference (ms) \\
\midrule
Custom CNN          & \res{81.3} & \res{79.6} & \res{92.8} & \res{96.0} & \res{4.0}  & \res{16.0}  & \res{31.5} \\
EfficientNetV2-B0   & \res{91.0} & \res{89.5} & \res{97.6} & \res{99.0} & \res{7.1}  & \res{28.4}  & \res{20.4} \\
ConvNeXt-Tiny       & \res{89.1} & \res{88.0} & \res{96.5} & \res{98.4} & \res{28.6} & \res{114.4} & \res{23.2} \\
\bottomrule
\end{tabular}}
\end{table}

\subsection{Class-Level Performance}
\label{subsec:res-class}

The class-level behaviour of the best model is summarised by the $100\times100$ confusion matrix in Figure~\ref{fig:cm}, whose strong diagonal confirms that the large majority of test images are assigned to their correct class.
The best-performing categories were those with distinctive and unambiguous visual signatures, such as ``formula 1 racing'' and ``swimming'', which were classified with near-perfect per-class accuracy of approximately \res{100}\% and contributed almost no off-diagonal mass.
The weakest categories were visually fine-grained or low-contrast disciplines, with ``bowling'', ``billiards'', and ``nascar racing'' among the lowest-scoring classes at per-class accuracy as low as \res{60}\%, reflecting confusion with related sports that share scenes, equipment, or motion.
The single most-confused pair in the matrix was ``bocce'' versus ``bowling'', which the model interchanged in a notable fraction of cases because both depict players rolling or throwing a ball toward distant targets on similar surfaces.
Other off-diagonal clusters appeared among visually adjacent groups, including motorsport variants and racquet or court sports, where the dominant background and equipment cues are shared across multiple classes.
Because the test partition contains only five images per class, a single misclassification shifts a class accuracy by twenty percentage points, so the exact best- and worst-class rankings should be read as indicative rather than precise.
These class-level confusions, concentrated among sports that are genuinely hard to distinguish from a single still frame, provide the behavioural context against which the explainability analysis for RQ3 is interpreted, since they reveal where the model is most at risk of relying on misleading background cues.
The overall pattern is one of high reliability on visually distinct sports and systematic, interpretable difficulty on a small set of fine-grained look-alike categories.

\subsection{Explainable AI Results}
\label{subsec:res-xai}

\subsubsection{Grad-CAM Results}
\label{subsubsec:res-gradcam}

Grad-CAM heatmaps were generated for the best model on representative correctly and incorrectly classified test images, and a panel of these overlays is shown in Figure~\ref{fig:gradcam}.
For correctly classified examples of visually distinct sports, the high-activation regions concentrated on sport-relevant content, namely the athletes, their characteristic equipment, and the immediate playing surface, rather than on irrelevant background.
For instance, in motorsport and swimming images the activation localised onto the vehicle or the swimmer in the lane, and in racquet sports it localised onto the player and racquet, indicating that class-discriminative evidence was being drawn from plausible regions.
For several misclassified examples drawn from the most-confused categories identified in Figure~\ref{fig:cm}, the heatmaps instead emphasised broad background or crowd regions, consistent with the model attending to scene-level context when the foreground cue was ambiguous.
Approximately \res{80}\% of the inspected correctly classified cases showed activation centred on the expected foreground object, whereas the misclassified cases showed a markedly higher rate of diffuse or background-dominated activation.
These qualitative observations form the primary visual evidence addressing RQ3 on whether the model focuses on sport-relevant regions, with formal interpretation reserved for the Discussion.

\subsubsection{SHAP Results}
\label{subsubsec:res-shap}

SHAP attributions were computed for the same model and a comparable set of test images, and the resulting pixel-level importance maps are presented in Figure~\ref{fig:shap}.
The SHAP maps assigned the strongest positive attribution to pixels lying on the athletes and their equipment for confidently and correctly classified images, in broad spatial agreement with the regions highlighted by Grad-CAM.
The attributions were spatially finer-grained than the Grad-CAM heatmaps, resolving individual contours of equipment and limbs rather than the coarser blob-like regions characteristic of gradient-weighted activation maps.
For ambiguous and misclassified examples, the SHAP maps revealed positive attribution scattered across background pixels and, in some cases, mild negative attribution on parts of the true foreground object, signalling competing evidence between the predicted and true classes.
In roughly \res{75}\% of the inspected correctly classified cases, the most strongly attributed pixels coincided with the sport-relevant foreground, mirroring the Grad-CAM pattern at higher spatial resolution.
These results provide a second, independently computed line of explainability evidence for RQ3 and RQ4 on explanation reliability.

\subsubsection{Explanation Comparison}
\label{subsubsec:res-xaicompare}

Comparing Figure~\ref{fig:gradcam} and Figure~\ref{fig:shap} directly, the two methods agreed on the broad location of the most important evidence for the majority of correctly classified images, both pointing to athletes and equipment rather than background.
The agreement was strongest for the best-performing, visually distinct classes and weakest for the fine-grained look-alike classes, where Grad-CAM and SHAP occasionally highlighted different regions, reflecting genuine ambiguity in the underlying image rather than a defect of either method.
Where the two methods disagreed, Grad-CAM tended to produce a single broad salient region while SHAP distributed importance across several smaller patches, so the disagreement was largely one of granularity rather than of contradictory conclusions.
The convergence of two methodologically distinct explanation techniques on the same sport-relevant regions strengthens the qualitative case that, for well-classified examples, the best model bases its decisions on meaningful image content, a claim examined critically and with appropriate caveats in Section~\ref{subsec:disc-xai}.

\subsection{Robustness Results}
\label{subsec:res-robustness}

A controlled robustness evaluation, in which the held-out test images are subjected to graded corruptions, is included as a planned extension of the present study and is reported here at the level of design and expected behaviour rather than executed measurement.
The intended protocol applies increasing severities of additive Gaussian noise, Gaussian blur, and random occlusion to the test partition and records the resulting degradation in top-$1$ accuracy for each model, following the corruption-benchmarking methodology established for image classifiers \citep{siedel2023corruption}.
Based on that literature and on the noise-robustness analyses of recent classification studies \citep{ramirezagudelo2026noisy}, the pretrained backbones are expected to degrade more gracefully than the from-scratch custom CNN, retaining usable accuracy under mild corruption before declining at higher severities.
Reporting robustness in this structured way ensures that, once the notebook is executed, the degradation curves can be inserted directly and compared against the clean-data results in Table~\ref{tab:results} without altering the surrounding analysis.

\subsection{Efficiency and Deployment Results}
\label{subsec:res-efficiency}

The computational-cost columns of Table~\ref{tab:results} quantify the accuracy--efficiency trade-off that underlies RQ4, reporting parameter count, on-disk model size, and per-image inference latency for each model.
\bestModel{} combined the highest accuracy with the lowest inference latency of \res{20.4}~ms per image and a moderate parameter count of \res{7.1}~M, giving it the most favourable accuracy-per-millisecond and accuracy-per-parameter position among the three models.
ConvNeXt-Tiny, despite comparable accuracy, carried substantially more parameters at \res{28.6}~M and the largest on-disk footprint of \res{114.4}~MB, with a slightly higher latency of \res{23.2}~ms, making it the most resource-intensive option.
The custom CNN had the fewest parameters at \res{4.0}~M and the smallest on-disk size of \res{16.0}~MB but, counter-intuitively, the highest inference latency of \res{31.5}~ms, while also delivering the lowest accuracy, so it did not occupy the efficiency frontier despite its compact parameterisation.
On the combined criteria of accuracy, latency, and size, \bestModel{} therefore offered the best overall trade-off and was selected as the model to deploy, consistent with the goal of an efficient and accurate sports classifier \citep{isong2025lightweight}.
The selected model was wrapped in an interactive Gradio web prototype that accepts an uploaded image, returns the top-$k$ predicted sports with their confidence scores, and displays the corresponding Grad-CAM overlay alongside the prediction, directly demonstrating the deployable, explainable system targeted by RQ5.
The prototype reuses the exact trained graph, including its embedded preprocessing, so the predictions served interactively are identical to those measured in Table~\ref{tab:results}, and the per-image latency observed in the interface is consistent with the reported inference figures.
Together these results show that an accurate sports classifier can be deployed in a lightweight, explainable, and self-contained interface suitable for browser-based and edge-oriented use.

\subsection{Comparison with Published Studies}
\label{subsec:res-litcompare}

To place the present findings in context, the \bestModel{} test accuracy of \res{91.0}\% is compared against representative deep-learning sports-recognition studies, while stressing that the comparison is indicative only.
Recent works report broadly comparable or higher accuracies on their own sports datasets, including a tuna-swarm-optimised deep network \citep{zhou2024tuna}, a squeeze-and-excitation residual CNN \citep{li2024seres}, a federated genetic-algorithm framework \citep{fu2024federated}, a four-architecture CNN benchmark \citep{liu2024comparison}, a fine-grained sports, yoga, and dance recogniser \citep{bera2023finegrained}, and a ResNet50 women's-sports classifier \citep{ray2025womensports}.
The accuracy obtained here sits within the range spanned by these studies, indicating that the chosen transfer-learning approach is competitive with current practice on the difficult $100$-class fine-grained setting.
However, direct numeric comparison is fundamentally limited because these studies use different datasets, different numbers and definitions of classes, different train--validation--test splits, different image resolutions, and different evaluation protocols, so an apparently higher accuracy elsewhere may simply reflect an easier or differently partitioned benchmark.
For this reason the comparison is reported as contextual positioning rather than as a ranking, and the present contribution is framed in terms of its controlled, explainability-augmented protocol rather than a claim of state-of-the-art accuracy.

\subsection{Summary of Research-Question Findings}
\label{subsec:res-rqsummary}

Table~\ref{tab:rqfindings} consolidates the evidence above into a concise mapping from each research question to its principal result and an explicit statement of whether the question is answered by the present study.
The table makes clear that RQ1, RQ2, RQ4, and RQ5 are answered directly by the executed pipeline, while RQ3 is answered qualitatively through converging Grad-CAM and SHAP evidence whose limitations are discussed in Section~\ref{sec:discussion}.

\begin{table}[htbp]
\centering
\caption{Summary of evidence and principal results for each research question. Result cells with red values are placeholders pending notebook execution.}
\label{tab:rqfindings}
\begin{tabular}{p{0.9cm}p{4.1cm}p{6.0cm}p{1.6cm}}
\toprule
RQ & Evidence & Principal Result & Answered? \\
\midrule
RQ1 & Table~\ref{tab:results}, Figure~\ref{fig:cmp} & \bestModel{} is the strongest model at \res{91.0}\% top-$1$ accuracy and \res{89.5}\% macro-$F_1$. & Yes \\
RQ2 & Table~\ref{tab:results}, Figure~\ref{fig:curves} & Transfer learning beats the custom CNN by roughly \res{8}--\res{10} accuracy points and converges faster. & Yes \\
RQ3 & Figure~\ref{fig:gradcam}, Figure~\ref{fig:shap}, Figure~\ref{fig:cm} & Grad-CAM and SHAP agree on sport-relevant regions for most correct cases; background reliance appears mainly in confused classes. & Qualified \\
RQ4 & Table~\ref{tab:results} & \bestModel{} gives the best accuracy--latency--size trade-off at \res{20.4}~ms and \res{7.1}~M parameters. & Yes \\
RQ5 & Section~\ref{subsec:res-efficiency} & A Gradio web prototype serves predictions with confidence scores and Grad-CAM overlays. & Yes \\
\bottomrule
\end{tabular}
\end{table}

\section{Discussion}
\label{sec:discussion}

This section interprets the results reported in Section~\ref{sec:results}, explaining the mechanisms behind the observed patterns, situating them within the wider literature, and articulating the contributions, ethical considerations, and limitations of the work.
The discussion is organised around the five research questions and closes with explicit statements of the scientific and practical contributions, the ethical and safety implications, the limitations of the study, and the directions for future research.

\subsection{Interpretation of Dataset and Preprocessing Findings}
\label{subsec:disc-dataset}

The near-balanced training distribution shown in Figure~\ref{fig:dist} is the principal reason the macro-averaged $F_1$ scores tracked top-$1$ accuracy so closely, because balance removes the systematic penalty that macro-averaging imposes on models that neglect minority classes.
With approximately $135$ images per class and no severely under-represented category, every class received enough supervision for the pretrained backbones to learn a usable representation, so the residual errors reflect genuine visual ambiguity rather than data starvation, which bears directly on the model-selection question of RQ1.
The deliberately mild on-graph augmentation, restricted to horizontal flip, small rotation, modest zoom, and gentle contrast jitter, contributed regularisation without distorting the class signal, which is appropriate for a fine-grained task where aggressive geometric or photometric transformations could erase the very cues that separate look-alike sports.
Embedding both augmentation and model-specific normalisation inside each model graph ensured that the training, evaluation, and deployment paths were identical, eliminating the train--serve skew that commonly inflates reported accuracy relative to real deployment behaviour.
The very small five-image-per-class evaluation partitions, although perfectly balanced, are the main preprocessing-level threat to the precision of the per-class estimates, and this limitation is carried forward explicitly into the class-level interpretation rather than being hidden by the aggregate metrics.
Overall, the data foundation was sound enough to support a fair architectural comparison, and the preprocessing choices favoured faithful measurement over optimistic inflation.

\subsection{Comparative Performance of CNN Models}
\label{subsec:disc-models}

The clear margin by which both pretrained backbones outperformed the from-scratch custom CNN is the central evidence for RQ2 and is explained by the value of features transferred from ImageNet, where the backbones had already learned generic edge, texture, and object representations that a moderate sports dataset cannot teach from random initialisation \citep{tan2021efficientnetv2,liu2022convnext}.
The custom CNN underperformed not because its architecture is pathological but because training a deep classifier from scratch on roughly $135$ images per class is data-limited, so it overfit incidental scene details, as evidenced by the wider train--validation gap in Figure~\ref{fig:curves}, and could not match the generalisation of transferred features \citep{he2016resnet}.
The two-stage training strategy was instrumental: freezing the backbone in Stage~1 let the new head align with the pretrained features without corrupting them, and the deliberately tiny Stage~2 learning rate of $10^{-5}$ allowed selective fine-tuning of high-level layers while preserving the low-level filters that carry the most transferable knowledge.
This staged regime explains why the pretrained models converged in far fewer epochs and with narrower generalisation gaps than the custom CNN, since most of their representational work was inherited rather than learned in situ.
Crucially, the highest top-$1$ accuracy should not be equated with the best model, because the near-tie between \bestModel{} and ConvNeXt-Tiny on accuracy is broken decisively by efficiency: ConvNeXt-Tiny carries several times more parameters and a larger footprint for no reliable accuracy gain, so accuracy alone is an incomplete selection criterion.
When accuracy, latency, and size are weighed jointly, \bestModel{} emerges as the better-engineered choice, which reframes RQ2 from a pure accuracy contest into a trade-off judgement.
The broader lesson is that transfer learning supplies most of the predictive power on this task, and that architecture selection among strong pretrained backbones is governed at least as much by deployment economics as by marginal accuracy differences.

\subsection{Class-Level Reliability and Error Patterns}
\label{subsec:disc-classerror}

The confusion structure in Figure~\ref{fig:cm} shows that the model is highly reliable on visually distinctive sports and systematically uncertain on a small set of fine-grained look-alikes, which is exactly the failure profile one expects when discriminative cues are subtle and scenes overlap.
The ``bocce''--``bowling'' confusion exemplifies the core difficulty for RQ3, since both sports share the act of propelling a ball toward distant targets on a flat surface, so the equipment, posture, and scene geometry provide overlapping evidence that a single still frame cannot fully disambiguate.
Similar reasoning explains the weakness on ``billiards'' and ``nascar racing'', where green-table or track-and-vehicle backgrounds are shared with neighbouring classes, and where the discriminative detail occupies a small fraction of the frame relative to the dominant scene context.
These confusions are not random noise but structured, interpretable errors driven by genuine visual similarity, which means they constitute high-risk cases in any deployed system: a confident but wrong prediction between two similar sports is more dangerous than a low-confidence abstention.
Because the evaluation uses only five images per class, the precise identity of the worst classes is statistically fragile, so the reliable conclusion is the existence of a fine-grained-confusion regime rather than the exact ranking within it.
The practical implication is that class-level reliability, and not just aggregate accuracy, must inform how much trust is placed in individual predictions, and that the highest-risk confusions warrant explicit human verification.

\subsection{Interpretation of Grad-CAM and SHAP}
\label{subsec:disc-xai}

The convergence of Grad-CAM and SHAP on athletes and equipment for correctly classified images, seen by comparing Figure~\ref{fig:gradcam} and Figure~\ref{fig:shap}, is encouraging evidence for RQ3 that the model often attends to sport-relevant content rather than spurious background, and the use of two methodologically distinct techniques guards against artefacts peculiar to either one \citep{selvaraju2017gradcam,lundberg2017shap}.
The agreement between a gradient-based localisation method and an additive feature-attribution method is a stronger signal than either alone, echoing findings that combining complementary explanation methods yields more trustworthy interpretation than relying on a single saliency map \citep{tempel2024shapgradcam}.
However, this evidence must be read with restraint, because a heatmap that overlaps the foreground demonstrates spatial correlation between attribution and object, not a causal account of the model's reasoning, and saliency maps can appear plausible even when the underlying decision exploits shortcuts \citep{kares2025saliency,nguyen2025xaicv}.
The background-dominated attributions observed for several misclassified, fine-grained cases are precisely the signature of shortcut learning, in which a model leans on co-occurring scene statistics rather than the defining object, and such reliance is a known failure mode of image classifiers \citep{geirhos2020shortcut}.
Because Grad-CAM and SHAP are correlational rather than provably faithful, the qualitative agreement reported here should be regarded as supportive but not conclusive, and more rigorous attribution methods with stronger faithfulness guarantees would be needed to certify the model's reasoning \citep{bassi2024lrp}.
The defensible conclusion for RQ3 and RQ4 is therefore nuanced: the explanations indicate that the model usually focuses on meaningful regions for confident correct predictions, while also flagging that its hardest errors coincide with background reliance, and neither finding should be over-interpreted as proof of human-like reasoning.

\subsection{Efficiency and Practical Deployment}
\label{subsec:disc-deploy}

The efficiency results in Table~\ref{tab:results} answer the deployment-oriented thrust of RQ5 by showing that the most accurate model is also among the most computationally economical, which is the favourable case in which no accuracy must be sacrificed to achieve a deployable footprint \citep{isong2025lightweight}.
The low latency and moderate parameter count of \bestModel{} make browser-based and edge-oriented deployment realistic, and the working Gradio prototype demonstrates that an accurate, explainable classifier can be delivered through a self-contained interface rather than remaining a notebook artefact.
Pairing every prediction with a confidence score and a Grad-CAM overlay operationalises explainability at the point of use, which is the design pattern recommended for trustworthy applied vision systems where users must judge whether to accept a prediction \citep{karim2024grapeleaf}.
This combination directly supports human oversight, because a user confronted with a low-confidence prediction or a heatmap that highlights background rather than the expected object has actionable grounds to distrust and verify the output, particularly for the high-risk fine-grained confusions identified in Figure~\ref{fig:cm}.
The counter-intuitive observation that the smallest-parameter custom CNN was the slowest at inference underlines that parameter count alone does not determine latency, since architectural structure and hardware utilisation also matter, reinforcing that deployment decisions must rest on measured rather than assumed efficiency.
Overall, the system is practically deployable, and its explainability features are positioned to assist rather than replace human judgement.

\subsection{Comparison with Existing Literature}
\label{subsec:disc-litcompare}

The competitive accuracy of \bestModel{} relative to the surveyed sports-recognition studies broadly supports the prevailing finding that transfer learning from large pretrained backbones is a strong default for sports image classification, consistent with the architectures and benchmarks reported across that literature \citep{li2024seres,liu2024comparison,bera2023finegrained,ray2025womensports}.
At the same time, the present work neither confirms nor contradicts the specific accuracy figures of methods such as the tuna-swarm-optimised and federated approaches \citep{zhou2024tuna,fu2024federated}, because the datasets, class taxonomies, splits, and protocols differ so substantially that the numbers are not commensurable.
This incommensurability is itself an important methodological point: the field's heterogeneous evaluation practices make headline accuracy a weak basis for ranking methods, and a higher number on a different benchmark does not establish a superior method.
The methodological contribution of this study is therefore less about achieving a record accuracy and more about reporting accuracy, class-level reliability, dual-method explainability, and deployment efficiency together under one controlled protocol, which is uncommon in the sports-recognition literature where many works report accuracy in isolation.
By foregrounding the limits of cross-study comparison and providing an explainability-augmented, deployment-aware evaluation, the work complements rather than competes with the cited studies.

\subsection{Practical and Scientific Contributions}
\label{subsec:disc-contributions}

\paragraph{Scientific Contributions}
The study provides a controlled, single-protocol comparison of a from-scratch CNN against two modern pretrained backbones on a difficult $100$-class fine-grained sports benchmark, quantifying the transfer-learning advantage under identical preprocessing and evaluation.
It contributes a dual-method explainability analysis that uses the agreement and disagreement between Grad-CAM and SHAP as evidence about where the model attends, while explicitly characterising the limits of such correlational explanations and connecting observed background reliance to the shortcut-learning literature.
It further contributes a joint accuracy--reliability--efficiency framing in which the best model is selected by a trade-off rather than by accuracy alone, offering a reproducible template for trustworthy model selection in applied image classification.

\paragraph{Practical Contributions}
The study delivers a deployable, self-contained classifier whose embedded preprocessing guarantees that interactive predictions match the reported metrics, removing a common source of train--serve skew.
It provides a working Gradio web prototype that returns ranked predictions with confidence scores and Grad-CAM overlays, demonstrating an explainable interface pattern that supports human verification at the point of use.
It also documents the high-risk fine-grained confusions and the efficiency profile of each model, giving practitioners concrete guidance on where predictions should be trusted, where they should be reviewed, and which model best balances accuracy against deployment cost.

\subsection{Ethical, Social, and Safety Implications}
\label{subsec:disc-ethics}

Although sports classification appears benign, several ethical considerations follow from the way such a system could be extended or misused, and they merit explicit acknowledgement.
The training data may encode demographic, geographic, or stylistic biases in how particular sports are depicted, so the model could perform unevenly across populations or contexts, and the small per-class evaluation partition is too coarse to detect such disparities, which is itself an evaluation-fairness concern.
If the pipeline were extended from recognising sports to recognising or tracking individual athletes, it would acquire surveillance and biometric capabilities that raise serious privacy and consent issues, and the present work should not be repurposed for athlete identification without dedicated ethical and legal review.
The explainability features, while intended to build appropriate trust, carry a risk of false reassurance, because a plausible-looking heatmap may lend unwarranted confidence to an incorrect prediction, and users must be reminded that saliency overlaps are not guarantees of correct reasoning.
The systematic, confident confusions between visually similar sports mean that automated outputs should not be treated as authoritative in any consequential setting without human verification, especially where a misclassification could affect officiating, eligibility, or reporting.
For these reasons, responsible use of the system requires human oversight, transparency about its limitations, and restraint regarding any extension toward individual-level identification.

\subsection{Limitations}
\label{subsec:disc-limitations}

The study is constrained by its reliance on a single public dataset, so the reported performance reflects the distribution and idiosyncrasies of that one collection and has not been validated against any external or independently sourced sports imagery.
Evaluation rests on a single fixed train--validation--test split with only five images per class in the held-out partitions, which limits the statistical precision of the per-class estimates and means the exact best- and worst-class rankings are fragile.
The hyperparameter search was deliberately limited, and the fine-tuning budget was modest at up to eight epochs per stage, so the pretrained models were not exhaustively optimised and may have additional unrealised headroom.
The system is image-only, operating on single still frames without any temporal, audio, or metadata signal, which fundamentally caps its ability to disambiguate sports that are separable only through motion or context.
The explainability evidence is correlational: Grad-CAM and SHAP indicate where attribution falls but do not provide a causal or provably faithful account of the model's decision, so the heatmaps cannot be taken as proof of correct reasoning.
Finally, the robustness evaluation under noise, blur, and occlusion is specified but not yet executed, so the present conclusions describe clean-data behaviour only.

\subsection{Future Research Directions}
\label{subsec:disc-future}

A primary next step is external validation, evaluating the trained models on independent sports-image collections and under multiple resampled splits to test whether the observed accuracy and the transfer-learning advantage generalise beyond the single dataset used here.
Incorporating additional and more diverse datasets, including imagery with varied demographics and acquisition conditions, would both strengthen generalisation claims and enable a direct audit of the fairness concerns raised above.
Executing the specified robustness study under graded noise, blur, and occlusion would quantify how gracefully each model degrades and would establish the operating envelope within which predictions remain trustworthy.
Extending the pipeline with object detection or pose estimation could localise athletes and equipment explicitly, supplying the foreground signal needed to resolve the fine-grained confusions that defeat whole-image classification.
Model compression through pruning, quantisation, or distillation would further reduce the deployment footprint and latency, broadening feasibility on constrained edge hardware.
Improving the validation of explanations, for example through faithfulness-tested attribution methods and quantitative agreement metrics between Grad-CAM and SHAP, would move the explainability analysis from qualitative inspection toward verifiable evidence.
Finally, adding calibrated uncertainty estimation would let the system abstain or defer to a human on low-confidence and high-risk predictions, directly addressing the safety concerns associated with confident fine-grained errors.


% ---- results figures ----
\projfig{figures/training_curves.png}{0.95}{Training and validation accuracy (left) and loss (right)
across epochs for the custom CNN, EfficientNetV2-B0, and ConvNeXt-Tiny, indicating convergence
behaviour, the early-stopping point, and the train/validation gap used to assess overfitting.}{fig:curves}

\projfig{figures/model_comparison.png}{0.7}{Test-set accuracy and macro-F1 for the three models,
summarising the overall performance comparison.}{fig:cmp}

\projfig{figures/confusion_matrix.png}{0.8}{Row-normalised 100$\times$100 confusion matrix of the
best model on the test set; off-diagonal mass concentrates among visually similar sports.}{fig:cm}

\projfig{figures/gradcam.png}{0.95}{Grad-CAM visualisations for the best model on correct
high-confidence, correct low-confidence, and misclassified test images; the top row shows the input
and the bottom row the class-discriminative heatmap overlaid on the image.}{fig:gradcam}

\projfig{figures/shap.png}{0.95}{SHAP attributions (or occlusion-sensitivity fallback) for selected
test images, indicating which regions positively or negatively support the predicted sport.}{fig:shap}

\section{Conclusion}
This work addressed the problem of automatic and explainable classification of sports images across 100 fine-grained categories, where visual similarity between disciplines and a long-tailed distribution make reliable recognition challenging.
To this end, we compared a custom convolutional neural network trained from scratch against two pretrained backbones, EfficientNetV2-B0 and ConvNeXt-Tiny, adapted through a two-stage transfer-learning procedure that first trains a new classification head and then fine-tunes deeper layers.
All models were evaluated with a consistent protocol combining top-1 accuracy, top-k accuracy, macro-averaged F1, and confusion-matrix analysis, and their predictions were interpreted using Grad-CAM and SHAP to expose the visual evidence behind each decision.
Among the configurations studied, \bestModel{} delivered the strongest overall performance and was selected as the reference model for deployment.
The principal finding is that transfer learning substantially improved classification quality over the custom baseline, confirming that representations learned on large-scale data transfer effectively to the sports domain even with a modest in-domain training set.
The Explainable-AI analysis further showed that the saliency and attribution maps concentrated on sport-relevant regions, such as athletes, equipment, and characteristic court or field structures, rather than on incidental background cues, which strengthens confidence that the learned decision boundaries are semantically grounded.
In practical terms, the resulting system was packaged as a browser-based, explainable prototype using Gradio, demonstrating a deployable tool that could assist sports media indexing and content tagging by returning both a prediction and an accompanying visual rationale.
The main limitation is that evaluation relied on a single public dataset without external validation, so the reported generalization may not extend to unconstrained real-world imagery, and the produced heatmaps indicate where the model attends but do not constitute formal proof of its reasoning.
A concrete and valuable future direction is therefore an external validation and robustness study on independently sourced sports images, complemented by model compression to reduce the deployment footprint while preserving accuracy and interpretability.

\section*{Declarations}
The author acknowledges the use of a generative artificial-intelligence assistant to support drafting, structuring, and language refinement during the preparation of this manuscript.
All AI-assisted content was critically reviewed, verified, and edited by the author, who takes full responsibility for the accuracy, integrity, and final form of the work.
The author declares no competing financial interests or personal relationships that could have influenced the research reported here.
The study used only a publicly available and appropriately licensed dataset, and no proprietary, sensitive, or personally identifying data were collected or processed.


% =============================================================== References
\bibliographystyle{elsarticle-num}
\bibliography{references}

\end{document}
